% Generated by tools/build_technical_appendix.py; edit the generator or source proofs.
\documentclass[11pt,a4paper]{article}
\usepackage{fontspec,amsmath,amssymb,mathtools}
\providecommand{\mathscr}[1]{\mathcal{#1}}
\DeclareFontFamily{U}{euf}{}
\DeclareFontShape{U}{euf}{m}{n}{<-> eufm10}{}
\usepackage[margin=22mm,headheight=14pt]{geometry}
\usepackage{graphicx,booktabs,longtable,array,calc,enumitem}
\usepackage{needspace,xurl,newunicodechar}
\usepackage[unicode,hidelinks]{hyperref}
\usepackage{bookmark,fancyhdr}
\defaultfontfeatures{Ligatures=TeX}
\usepackage{xeCJK}
\setmainfont{Times New Roman}
\setsansfont{Arial}
\setCJKmainfont[Path=../fonts/,
  BoldFont=NotoSansSC-SemiBold.ttf,
  ItalicFont=NotoSerifSC-Regular.ttf]{NotoSerifSC-Regular.ttf}
\setCJKsansfont[Path=../fonts/,
  BoldFont=NotoSansSC-SemiBold.ttf]{NotoSansSC-SemiBold.ttf}
\setCJKmonofont[Path=../fonts/]{NotoSansSC-SemiBold.ttf}
\xeCJKsetup{PunctStyle=quanjiao,CJKmath=true}
\renewcommand{\figurename}{图}
\XeTeXlinebreaklocale "zh"
\XeTeXlinebreakskip=0pt plus 1pt
\setmonofont{Consolas}[Scale=0.82]
\newfontfamily\extendedlatin{lmroman10-regular.otf}
\newunicodechar{ő}{{\extendedlatin\char"0151}}
\newunicodechar{ű}{{\extendedlatin\char"0171}}
\setlength{\emergencystretch}{4em}
\setlength{\parindent}{1.5em}
\setlength{\parskip}{2.5pt plus 1pt}
\linespread{1.10}
\allowdisplaybreaks[3]
\raggedbottom
\clubpenalty=10000
\widowpenalty=10000
\displaywidowpenalty=10000
\setcounter{secnumdepth}{-1}
\providecommand{\tightlist}{\setlength{\itemsep}{0pt}\setlength{\parskip}{0pt}}
\newcommand{\passthrough}[1]{#1}
\setlist{itemsep=2pt,parsep=0pt}
\makeatletter
\renewcommand\section{\@startsection{section}{1}{\z@}{-2.8ex plus -1ex minus -.2ex}{1.1ex plus .2ex}{\normalfont\large\bfseries}}
\renewcommand\subsection{\@startsection{subsection}{2}{\z@}{-2.1ex plus -.8ex minus -.2ex}{.8ex plus .2ex}{\normalfont\normalsize\bfseries}}
\renewcommand\subsubsection{\@startsection{subsubsection}{3}{\z@}{-1.8ex plus -.5ex minus -.2ex}{.6ex plus .2ex}{\normalfont\normalsize\bfseries}}
\makeatother
\pagestyle{fancy}
\fancyhf{}
\fancyhead[C]{\small 竞赛图 Hamilton 路径计数技术附录}
\fancyfoot[C]{\small\thepage}
\renewcommand{\headrulewidth}{0pt}
\hypersetup{pdftitle={竞赛图 Hamilton 路径计数的技术附录},pdfauthor={刘星辰、叶祥宇}}
\begin{document}
\thispagestyle{plain}
\begin{center}
{\Large\bfseries 竞赛图 Hamilton 路径计数的技术附录\par}
\vspace{9pt}
{\normalsize 刘星辰、叶祥宇\par}
\vspace{5pt}
{\small 2026年11月11日\par}
\end{center}
\Needspace{12\baselineskip}
\section{范围与依赖导航}\label{ux8303ux56f4ux4e0eux4f9dux8d56ux5bfcux822a}

本附录集中收录主文《竞赛图 Hamilton 路径的定量计数与度稳定性》实际使用的较长分析证明。附录完整呈现所需证明，保留章节号、定理号、式号和引文号，以便与主文互相引用。本附录与主文共同构成所需证明；详细稿的其余部分作为可选展开阅读，不是这些证明的额外前提。

永久式近似在这里作为分析工具使用，附录给出固定元素界和谱隙参数下的独立定量证明。本文不把行列式主项称为新结果，也不在此提出关于适用矩阵范围的更广新颖性主张。第 3 节的局部文献比较与主文第 3 节的贡献定位一致。

\begin{itemize}
\tightlist
\item
  \hyperref[technical-permanent]{第 3 节}完整证明定理 3.1，包括高斯配对、压缩后的非高斯核心、线性系数窗口、阶乘恢复、系数尾部及最终一致误差界。
\item
  \hyperref[technical-scaling]{第 4 节}完整证明局部对角缩放及其代价、高斯删除与中心化、配对度差预处理、非主子阵的精确恢复恒等式及小度差近似。引理 4.4 使用定理 3.1 和引理 4.1--4.3；引理 4.5 使用相同的缩放与删除估计，但不作配对预缩放。
\item
  \hyperref[technical-spectral-cap]{引理 6.1}保留完整陈述与证明，以补齐 A.1 对传递矩阵行列式多项式的依赖。其中的特征值计算也给出第 4.3 节圆环例子所用的频率。
\item
  \hyperref[technical-algebra]{A.1--A.3}给出主文谱论证所需的传递矩阵行列式恒等式、交替切换同构、矩阵幂的对角元相等及四次矩等号代数。第 4.3 节的圆环例子直接使用 A.2 和引理 6.1。
\end{itemize}

未特别注明的章节和引理引用均指本附录；对主文的引用明确标出。所需记号约定在下文重述。标准外部工具保留原引文。附录完整保留原参考文献表以维持引文编号，不收录原作者说明和个人致谢。

\Needspace{12\baselineskip}
\section{记号}\label{ux8bb0ux53f7}

最大度差与主文统一记为 \(\Delta_s=\|S\mathbf1\|_\infty\)。A.2--A.3 中的 \(d\) 表示矩阵或竞赛图阶数；第 3 节的局部块大小和度数指标也保留 \(d\)。标量分析控制函数记为

\[
\mathcal G_{C,q}(r)=\exp\!\left(\frac{C^2r^2}{2(1-q^2r^2)}\right),\qquad rq<1.
\]

它既不同于第 4 节的度差乘积 \(\Gamma=\prod_i(1-a_i^2)\)，也不同于高斯矩阵泛函 \(\mathcal G(M)=\det(I-MM^{\mathsf T})^{-1/2}\)。第 3 节的生成函数 \(G(t)=\det(I-t^2BB^{\mathsf T})^{-1/2}\) 也是局部记号。该节的 \(C,q\) 是固定分析参数；第 4 节中 \(C=2A/(n-1)\) 是矩阵，A.2 中 \(q=(d-1)/2\) 是整数。

对竞赛图邻接矩阵 \(A\)，记

\[
S=A-A^{\mathsf T},\qquad 2A=J-I+S,\qquad s=S\mathbf1,
\qquad s_i=2d_i^+-(n-1),\qquad \sum_i s_i=0,
\] \[
\mathcal E=\|s\|_2^2,\qquad
\Delta_s=\|s\|_\infty,\qquad
a=\frac{s}{n-1},\qquad \tau=\|a\|_2^2\quad(n\ge2).
\]

这里 \(I_p\) 是单位矩阵，\(J_p=\mathbf1_p\mathbf1_p^{\mathsf T}\)，\(P_p=J_p/p\)，\(\Pi_p=I_p-P_p\)。子矩阵行列均按原顶点标号递增排列，行集与列集可以不同。第 4 节中 \(n\) 或 \(N\) 是完整阶数，\(m=n-t\) 是保留阶数，\(p\) 是一般矩阵阶数。空和为零，空积为一，空矩阵的行列式与永久式均为一。\(\|\cdot\|_{\rm op}\) 表示 Euclidean 算子范数，\(\|\cdot\|_F\) 或 \(\|\cdot\|_{\mathrm F}\) 表示 Frobenius 范数。

若 \(S\) 的非零频率对为 \(\pm i\lambda_j\)，记

\[
x_j=\lambda_j^2/n^2,\qquad
D_n(S)=\det(I-S^{\mathsf T}S/n^2)^{-1/2}
=\prod_j(1-x_j)^{-1},
\] \[
\rho_n(S)=D_n(S)\det(I+S/n)
=\prod_j\frac{1+x_j}{1-x_j}.
\]

奇异值成对给出 \(\sum_jx_j=(n-1)/(2n)<1/2\)，所以这些行列式因子为正；平方根均取正根。\(T_n^0\)、\(T_n\)、\(T_d\) 表示对应阶数的传递符号矩阵，即严格上三角元素全为 \(+1\) 的反对称矩阵。奇阶圆环图的阶数为 \(d=2q+1\)，其边 \(i\to j\) 当且仅当 \(j-i\pmod d\in\{1,\ldots,q\}\)。对数均为自然对数。一致误差常数可以依赖明确固定的分析参数或删除窗口参数，但不依赖具体的容许矩阵、竞赛图或删除集合。

\phantomsection\label{technical-permanent}

\Needspace{12\baselineskip}
\section{3. 永久式的一致渐近估计}\label{ux6c38ux4e45ux5f0fux7684ux4e00ux81f4ux6e10ux8fd1ux4f30ux8ba1}

本节证明一个关于永久式的统一估计。将行列和均为一的矩阵写成 \((J_n+E)/n\)，其中 \(E\mathbf1=E^{\mathsf T}\mathbf1=0\)。在元素有界且 \(\|E/n\|_{\mathrm{op}}\le q<1\) 的条件下，永久式的主项由一个行列式给出。第 4 节将把这一估计用于缩放后的竞赛图子矩阵。

McCullagh {[}5, §3{]} 已在适度偏离假设下给出具有 \(O(n^{-1})\) 误差的行列式近似。对非负输入，本文的固定元素界和谱隙假设蕴含弱有界性及所需谱隙，其行列式因子恰为 \(\det(I-BB^{\mathsf T})^{-1/2}\)；主文第 3 节详细核对这一对应。固定矩阵类上的一致性本身也不能构成相对于逐个容许序列定理的新颖性区别。这里保留独立证明，以逐步跟踪压缩核心、系数窗口和尾部中的常数，供删除与缩放之后使用，并不声称修复 {[}5{]} 中已确认的缺口。Li {[}8, Theorem 2.1{]} 假设中心扰动的最大绝对行和或列和为 \(o(n)\)；这一较窄条件不覆盖这里使用的稠密核，但这一比较不构成相对于 {[}5{]} 或 {[}12{]} 的新颖性主张。

压缩核心证明独立形成于此前的路径计数项目 {[}14{]}。有符号输入不构成优先性区别：McCullagh {[}5, §4.1{]} 已讨论复扰动。这里保留实有符号版本及其证明；竞赛图应用实际使用非负输入。例如

\[
B=\frac17\begin{pmatrix}1&2&-3\\-2&1&1\\1&-3&2\end{pmatrix},\qquad E=3B
\]

满足 \(B\mathbf1=B^{\mathsf T}\mathbf1=0\)、\(\|B\|_{\rm op}=5/7<1\)，而 \(J_3+E\) 含负元素。对任意 \(m\ge1\)，令 \(E_m=E\otimes J_m\)、\(n=3m\)，则 \(E_m\) 双中心化，条目绝对值至多 \(9/7\)，且 \(\|E_m/n\|_{\rm op}=5/7\)，而 \(J_n+E_m\) 仍含负元素。因此同一组固定参数在任意大的阶数上都容许有符号输入；这个例子区分两类假设，不承担渐近结论的证明。第 3.6 节另解释从 {[}12{]} 的非负版本证明提取定量误差的途径；这里不主张该方法原则上不能处理有符号输入。

\Needspace{12\baselineskip}
\subsection{3.1 定理与证明思路}\label{ux5b9aux7406ux4e0eux8bc1ux660eux601dux8def}

\textbf{定理 3.1（永久式的一致渐近估计）。}

固定 \(0\le C<\infty\) 和 \(0\le q<1\)。设 \(n\ge1\)，\(E\in\mathbb R^{n\times n}\) 满足

\[
E\mathbf1=E^{\mathsf T}\mathbf1=0,\qquad
\max_{i,j}|E_{ij}|\le C,\qquad
\|E/n\|_{\mathrm{op}}\le q.
\]

令 \(B=E/n\)。对所有满足上述条件的矩阵，一致地有

\begin{equation}
\frac{\operatorname{per}(J_n+E)}{n!}
=\det(I-BB^{\mathsf T})^{-1/2}+O_{C,q}(n^{-1}).
\tag{3.1}
\label{eq-uniform-permanent}
\end{equation}

\hyperref[eq-uniform-permanent]{式 (3.1)}也可写成行列式因子乘以相对误差 \(1+O_{C,q}(n^{-1})\)。下标 \(C,q\) 表示误差常数和起始阶数可依赖这两个预先固定的参数，但须同时适用于所有满足条件的矩阵。

本节所用行列式平方根，在有关实区间上均取正平方根。假设给出 \(\|B\|_{\mathrm F}^2\le C^2\)。因此，行列式因子不小于一，并且被只依赖 \(C,q\) 的常数一致控制。

对行、列指标的互异性分别作容斥，永久式的每个系数可表示为二部多重图之和。中心化消去一度顶点；二度分量对应高斯配对，其总和给出行列式主项。其余分量按超额（边数减顶点数）分类，每单位超额带来一个 \(n^{-1}\) 因子。我们先估计次数不超过 \(\alpha n\) 的系数，再以复分析方法控制其余系数。

证明中使用高斯生成级数 \(G(t)\)。令

\[
f(t)=\frac{\operatorname{per}(J_n+tE)}{n!},
\qquad G(t)=\det(I-t^2BB^{\mathsf T})^{-1/2}.
\]

选取 \(1<\sigma<R<1/q\)；当 \(q=0\) 时不要求上界。例如可取

\[
R=\frac{3+q}{2(1+q)},\qquad \sigma=\frac{1+R}{2}.
\]

使用前面定义的标量分析控制函数 \(\mathcal G_{C,q}\)：

\[
\mathcal G_{C,q}(r)=\exp\left(\frac{C^2r^2}{2(1-q^2r^2)}\right),
\qquad r\ge0,\ rq<1.
\]

本节在 \(q=0\) 时均省去涉及 \(1/q\) 的上界限制。所有常数仅依赖 \(C,q\) 及所选半径。第 3.6 节将它们合并为显式误差界。

对上述半径，

\[
R-1=\frac{1-q}{2(1+q)}>0,\qquad
1-Rq=\frac{2-q-q^2}{2(1+q)}=\frac{(1-q)(q+2)}{2(1+q)}>0,
\]

\[
\sigma-1=(R-1)/2>0,\qquad R-\sigma=(R-1)/2>0.
\]

以下以 \(s_i\) 表示 \(B\) 的奇异值。则

\[
\sum_i s_i^2=\sum_{i,j}B_{ij}^2\le n^2(C/n)^2=C^2,\qquad
0\le s_i\le q.
\]

由 \(-\log(1-u)=\int_0^u(1-v)^{-1}\,dv\le u/(1-u)\)，

\[
\begin{aligned}
\log G(r)&=-\frac12\sum_i\log(1-r^2s_i^2)\\
&\le\frac12\sum_i\frac{r^2s_i^2}{1-r^2s_i^2}
\le\frac{r^2\sum_i s_i^2}{2(1-r^2q^2)}
\le\frac{C^2r^2}{2(1-r^2q^2)}.
\end{aligned}
\]

因此 \(1\le G(1)\le\mathcal G_{C,q}(1)\)。若 \(|f(1)-G(1)|\le K/n\)，则

\[
\left|\frac{f(1)}{G(1)}-1\right|
=\frac{|f(1)-G(1)|}{G(1)}\le K/n,
\]

因此绝对误差界也给出同阶的相对误差界。

\Needspace{12\baselineskip}
\subsection{3.2 系数与高斯配对}\label{ux7cfbux6570ux4e0eux9ad8ux65afux914dux5bf9}

\textbf{定义 3.1（永久式系数与重标度）。}

写 \(f(t)=\sum_{k=0}^n a_kt^k\)。按永久式中来自 \(E\) 的元素展开，得到

\[
a_k=\frac1{(n)_k}
\sum_{\substack{I,J\subseteq[n]\\|I|=|J|=k}}
\operatorname{per}E[I,J],
\]

这里 \(0\le k\le n\)，所以分母 \((n)_k\) 非零。对 \(k>n\) 直接令 \(a_k=0\)，不使用含 \(1/(n)_k\) 的表达式。定义

\[
F_k=\frac{(n)_k}{n^k}a_k
=\frac1{k!}
\sum_{\substack{i_1,\ldots,i_k\ \mathrm{distinct}\\
j_1,\ldots,j_k\ \mathrm{distinct}}}
\prod_{\ell=1}^kB_{i_\ell j_\ell},
\qquad
F(t)=\sum_{k=0}^nF_kt^k.
\]

当 \(k>n\) 时，右侧为空和，按约定取零；以下形式恒等式也采用这一约定。

\textbf{定义 3.2（集合分划与权重）。}

对任意有限集合 \(E\)，记 \(\mathfrak P(E)\) 为将 \(E\) 拆成互不相交非空块的全部集合分划，并记 \(\mathfrak P_k=\mathfrak P([k])\)。``\(k\) 个坐标互异''的示性函数具有标准的分划格容斥展开。{[}5, Section 8{]} 的永久式计算也使用成对分划展开。其中，大小为 \(d\) 的块带 Möbius 权重 \((-1)^{d-1}(d-1)!\)；整个分划的权重在此定义为 \(\mu(\pi)=\prod_{D\in\pi}(-1)^{|D|-1}(|D|-1)!\)。分别对 \(F_k\) 的行坐标和列坐标应用这一展开。每对分划对应一个二部多重图：边带互异标号 \(1,\ldots,k\)，两侧的顶点分别为行、列分划的块，每个大小为 \(d\) 的块带上述权重。

每个顶点的数值标签独立地在 \([n]\) 中求和；不同顶点的数值标签允许相同。如果某个行顶点的度数为一，对其标签求和得到 \(B\) 的一个列和，因而为零。度数为一的列顶点同样给出零行和。因此，只有所有顶点度数至少为二的图留下贡献。

例如，在二次和三次项中，保留下来的行、列分划各自都只有一个块。两侧 Möbius 权重的乘积分别为 \(1\) 和 \(4\)，再除以 \(k!\)，得到

\[
F_2=\frac12\sum_{i,j}B_{ij}^2,
\qquad
F_3=\frac23\sum_{i,j}B_{ij}^3,
\qquad |F_3|\le\frac{2C^3}{3n}.
\]

其中二次项是高斯主项的一部分，三次项则具有 \(O(n^{-1})\) 的界。

所有顶点度数均为二的连通分量称为纯二度分量，其总指数生成函数为 \(G(t)\)。可用两个独立标准实高斯向量 \(X,Y\) 直接验证：由 Wick 公式，

\[
\mathbb E\exp(tX^{\mathsf T}BY)
\]

中 \(t^k\) 的系数恰等于边标号两侧的配对分划之和，而这些正是所有顶点均为二度的图。先对条件高斯分布积分，再将 \(BB^{\mathsf T}\) 对角化，得到

\[
\mathbb E\exp(tX^{\mathsf T}BY)
=\mathbb E\exp\!\left(\frac{t^2}{2}X^{\mathsf T}BB^{\mathsf T}X\right)
=\det(I-t^2BB^{\mathsf T})^{-1/2}.
\]

该解析恒等式在 \(|t|\|B\|_{\mathrm{op}}<1\) 时成立，并确定全部形式系数。特别地，\(G\) 的系数非负，且

\[
G(r)\le\mathcal G_{C,q}(r)\qquad(r\ge0,\ rq<1),
\]

因为 \(-\log(1-x)\le x/(1-x)\)，且 \(\sum_i s_i^2=\|B\|_{\mathrm F}^2\le C^2\)；这里 \(s_i\) 为 \(B\) 的奇异值。

从每个留下贡献的图中删除全部纯二度分量，称剩余部分为核心。核心可以不连通。如果核心有 \(h\) 条边、\(v\) 个顶点，定义其超额为 \(j=h-v\)。非空核心满足 \(j\ge1\)。带边标号的分量分解给出逐系数恒等式

\[
F(t)=G(t)\left(1+\sum_{j\ge1}C_j(t)\right),
\]

其中 \(C_j\) 是超额为 \(j\) 的核心生成级数，系数包括 Möbius 符号和矩阵收缩。该等式在形式级数意义下成立：固定次数 \(k\) 只涉及有限多个 \(j<k\)，无须假设关于 \(j\) 的级数在非零 \(t\) 处收敛。由 \(F_k\) 的互异指标定义，所有高于 \(n\) 次的系数均为零。

\textbf{定义 3.3（核心级数）。}

对 \(\pi,\sigma\in\mathfrak P_h\)，以 \(D_\pi(a)\)、\(D_\sigma(a)\) 表示含边标号 \(a\) 的两个分划块。令行块标签为 \(x_D\in[n]\)，列块标签为 \(y_E\in[n]\)，定义收缩

\[
\mathcal W_B(\pi,\sigma)
=\sum_{x:\pi\to[n]}\sum_{y:\sigma\to[n]}
\prod_{a=1}^hB_{x_{D_\pi(a)},y_{D_\sigma(a)}}.
\]

两个求和遍历全部函数，不要求单射。定义 \(\mathcal K_{j,h}\) 为满足下列条件的分划对：每个块大小至少为二；对应二部多重图的每个连通分量至少含一个大小不小于三的块；\(h-|\pi|-|\sigma|=j\)。其中 \(|\pi|\) 是块数，不是原集合大小。连通分量指能通过边依次相接到达的顶点组成的极大集合；度数按入射半边数计，重边分别计数，压缩后自环向其顶点贡献二度。核心级数定义为

\[
C_j(t)=\sum_{h=0}^{\infty}c_{j,h}t^h,\qquad
c_{j,h}=\frac1{h!}\sum_{(\pi,\sigma)\in\mathcal K_{j,h}}
\mu(\pi)\mu(\sigma)\mathcal W_B(\pi,\sigma).
\]

对固定 \(h\)，上述和为有限和。非空核心两侧至少各含一个块，故 \(j\le h-2\)，从而 \(h\le j+1\) 时 \(c_{j,h}=0\)。以下也使用较弱的条件 \(j<h\)。空图的贡献为 \(F_0=G(0)=1\)，不计入 \(C_j\)。

记 \(g_k=[t^k]G(t)\)，\(k\ge0\)，其中 \([t^k]\) 表示系数提取。于是 \(g_0=1\)，且奇数 \(k\) 时 \(g_k=0\)。引理 3.2 将证明加权范数 \(\|C_j\|_R=\sum_h|c_{j,h}|R^h\) 有限。

固定从 \(E\) 中选出的行集 \(I\)、像集 \(J\)，大小均为 \(k\)。在这些位置上选择双射的总权重是 \(\operatorname{per}E[I,J]\)；其余 \(n-k\) 行可以任意双射到其余列，总共有 \((n-k)!\) 种。因此

\[
[t^k]\operatorname{per}(J+tE)
=(n-k)!\sum_{|I|=|J|=k}\operatorname{per}E[I,J],
\]

\[
a_k=\frac{(n-k)!}{n!}\sum_{|I|=|J|=k}\operatorname{per}E[I,J]
=\frac1{(n)_k}\sum_{|I|=|J|=k}\operatorname{per}E[I,J].
\]

给定行集、列集及一个双射，将这 \(k\) 条匹配边依次编为 \(1,\ldots,k\) 有 \(k!\) 种方式；故

\[
\sum_{\boldsymbol i,\boldsymbol j\ \mathrm{distinct}}\prod_{h=1}^kB_{i_hj_h}
=k!\sum_{|I|=|J|=k}\operatorname{per}B[I,J]
=\frac{k!}{n^k}\sum_{|I|=|J|=k}\operatorname{per}E[I,J].
\]

对任一条件 \(\mathcal A\)，\(\mathbf1\{\mathcal A\}\) 在条件成立时为一，否则为零；它是标量示性函数，和全一向量 \(\mathbf1_p\) 不同。对分划 \(\pi\) 定义

\[
\mu(\pi)=\prod_{D\in\pi}(-1)^{|D|-1}(|D|-1)!,\quad
\mathbf1_\pi(\boldsymbol i)=\prod_{D\in\pi}\mathbf1\{i_h\text{ 在 }D\text{ 上相等}\}.
\]

则

\[
\mathbf1\{i_1,\ldots,i_k\text{ 互异}\}
=\sum_{\pi\in\mathfrak P_k}\mu(\pi)\mathbf1_\pi(\boldsymbol i).
\]

验证方式如下。设给定坐标的相等类为 \(D_1,\ldots,D_r\)，则右端分解为

\[
\prod_{a=1}^r\left(\sum_{\pi\in\mathfrak P(D_a)}\mu(\pi)\right).
\]

置换符号 \(\operatorname{sgn}(\pi)=(-1)^{\#\{(i,j):i<j,\ \pi(i)>\pi(j)\}}\)，即逆序对数为偶数时取一、奇数时取负一。大小为 \(d\) 的块可以放成一个有向循环的方式为 \((d-1)!\)，循环的置换符号为 \((-1)^{d-1}\)。因此括号内等于该集合上全部置换的符号和：大小一时为一，大小至少二时用乘一个固定换位的双射使正负项成对抵消，值为零。

两侧代入得

\[
F_k=\frac1{k!}\sum_{\pi,\sigma\in\mathfrak P_k}\mu(\pi)\mu(\sigma)
\sum_{\boldsymbol i,\boldsymbol j}\mathbf1_\pi(\boldsymbol i)\mathbf1_\sigma(\boldsymbol j)
\prod_{h=1}^kB_{i_hj_h}.
\]

若 \(\pi\) 有单元素块 \(\{h\}\)，对仅出现一次的标签 \(i_h\) 求和便得到 \(\sum_iB_{ij_h}=0\)。列侧同理。对于 \(k=2\)，只剩一对二元素块，权重 \((-1)(-1)=1\)；对于 \(k=3\)，只剩一对三元素块，权重 \(2\cdot2=4\)。故

\[
F_2=\frac1{2!}\sum_{i,j}B_{ij}^2,\quad
F_3=\frac4{3!}\sum_{i,j}B_{ij}^3,\quad
|F_3|\le\frac23n^2(C/n)^3=\frac{2C^3}{3n}.
\]

令 \(K=BB^{\mathsf T}\)、\(Q=B^{\mathsf T}B\)。四次项为

\[
\begin{aligned}
F_4={}&\frac18(\operatorname{tr}K)^2+\frac14\operatorname{tr}(K^2)\\
&-\frac34\sum_i\left(\sum_jB_{ij}^2\right)^2
-\frac34\sum_j\left(\sum_iB_{ij}^2\right)^2
+\frac32\sum_{i,j}B_{ij}^4.
\end{aligned}
\]

前三类分划配对分别为：两侧都是 \(2+2\)（\(3\cdot3=9\) 对）；一侧为 \(4\)、另一侧为 \(2+2\)（每个方向三对，块权重 \(-6\)）；两侧都是 \(4\)（权重 \(36\)）。除以 \(4!=24\) 后得到上述系数。两侧配对相同的三对给 \(3/24=1/8\)；配对不同的六对给 \(6/24=1/4\)。其中 \(g_4\) 为高斯主项，其余为非高斯修正。

标准实高斯变量 \(X\) 的密度为 \((2\pi)^{-1/2}e^{-x^2/2}\)。先设 \(u\in\mathbb R\)，配方得

\[
\mathbb Ee^{uX}=\frac1{\sqrt{2\pi}}\int_{\mathbb R}e^{-x^2/2+ux}\,dx
=e^{u^2/2}\frac1{\sqrt{2\pi}}\int_{\mathbb R}e^{-(x-u)^2/2}\,dx=e^{u^2/2}.
\]

对 \(v\in\mathbb R^n\)，独立标准向量的多元母函数为 \(\mathbb Ee^{v^{\mathsf T}X}=e^{\|v\|^2/2}\)。对坐标求导得到

\[
\mathbb E\prod_{h=1}^{2r}X_{i_h}
=\sum_{\mathcal P\text{ 为配对}}\prod_{\{a,b\}\in\mathcal P}\mathbf1\{i_a=i_b\},
\qquad \mathbb E\prod_{h=1}^{2r+1}X_{i_h}=0.
\]

实高斯矩可由一维积分计算。标准实高斯向量的坐标相互独立，密度均为 \((2\pi)^{-1/2}e^{-x^2/2}\)。令

\[
m_r=\frac1{\sqrt{2\pi}}\int_{\mathbb R}x^re^{-x^2/2}\,dx.
\]

归一化给出 \(m_0=1\)。奇函数对称积分给出 \(m_{2r+1}=0\)。对 \(r\ge1\) 分部积分：

\[
\begin{aligned}
m_{2r}
&=-\frac1{\sqrt{2\pi}}\int_{\mathbb R}x^{2r-1}\,d(e^{-x^2/2})\\
&=\frac{2r-1}{\sqrt{2\pi}}\int_{\mathbb R}x^{2r-2}e^{-x^2/2}\,dx\\
&=(2r-1)m_{2r-2}
=(2r-1)(2r-3)\cdots1.
\end{aligned}
\]

这里边界项为零，因为 \(|x|^ae^{-x^2/2}\to0\) 对每个固定整数 \(a\ge0\) 成立。记 \((-1)!!=1\)，于是统一写成 \(m_{2r}=(2r-1)!!\)，包括 \(r=0\)。

对指标列 \(i_1,\ldots,i_k\)，定义重数 \(c_a=\#\{h:i_h=a\}\)。独立性给出有限乘积

\[
\mathbb E\prod_{h=1}^kX_{i_h}=\prod_{a=1}^n m_{c_a}.
\]

任一 \(c_a\) 为奇数时右侧为零。若全部 \(c_a\) 为偶数，所有允许配对必须把相同指标的位置配在一起。\(c_a\) 个位置的配对数满足

\[
p_0=1,\qquad p_{2r}=(2r-1)p_{2r-2}
=(2r-1)!!.
\]

所以允许配对的总数是 \(\prod_a(c_a-1)!!\)，与前面算得的矩相同。定义 Kronecker 符号 \(\delta_{ab}=\mathbf1\{a=b\}\)，便得到

\[
\mathbb E\prod_{h=1}^{2r}X_{i_h}
=\sum_{\mathcal P}\prod_{\{a,b\}\in\mathcal P}\delta_{i_ai_b}.
\]

这里 \(\mathcal P\) 遍历 \([2r]\) 的配对；\(r=0\) 时空配对贡献一，奇数阶矩为零。

先取实数 \(t\) 满足 \(|t|q<1\)，对 \(Y\) 积分，然后对 \(BB^{\mathsf T}\) 作正交对角化。标准向量密度为 \((2\pi)^{-n/2}e^{-\|x\|_2^2/2}\)；正交变换保持 \(\|x\|_2\) 和 Jacobian 绝对值一，故变换后的密度相同且仍为各坐标密度的乘积。因此：

\[
\begin{aligned}
\mathbb Ee^{tX^{\mathsf T}BY}
&=\mathbb E_Xe^{t^2\|B^{\mathsf T}X\|^2/2}\\
&=\prod_i\frac1{\sqrt{2\pi}}\int_{\mathbb R}
e^{-(1-t^2s_i^2)x^2/2}\,dx\\
&=\prod_i(1-t^2s_i^2)^{-1/2}=G(t).
\end{aligned}
\]

为将等式推广到复圆盘，令 \(U=X^{\mathsf T}BY\)，并取实数 \(r\) 满足 \(|t|<r<1/q\)。当 \(q=0\) 时 \(B=0\)，结论直接成立。对实数 \(u\)，有

\[
e^{r|u|}\le e^{ru}+e^{-ru}.
\]

由于 \(Y\) 与 \(-Y\) 同分布，\(U\) 与 \(-U\) 同分布，所以

\[
\mathbb E e^{r|U|}
\le\mathbb E e^{rU}+\mathbb E e^{-rU}
=2G(r)<\infty.
\]

由指数级数各项非负，对 \(k\ge0\) 有

\[
\frac{r^k|U|^k}{k!}\le e^{r|U|}
\quad\Longrightarrow\quad
\mathbb E|U|^k\le\frac{k!}{r^k}\mathbb Ee^{r|U|}
\le\frac{2k!}{r^k}G(r)<\infty.
\]

因此对每个整数 \(N\ge0\) 都可以对有限和逐项取期望：

\[
\sum_{k=0}^{N}\frac{r^k}{k!}\mathbb E|U|^k
=\mathbb E\sum_{k=0}^{N}\frac{r^k|U|^k}{k!}
\le\mathbb Ee^{r|U|}\le2G(r).
\]

再令 \(s=|t|/r<1\)。逐点控制指数级数的余项：

\[
\left|e^{tU}-\sum_{k=0}^{N}\frac{t^kU^k}{k!}\right|
\le\sum_{k>N}s^k\frac{r^k|U|^k}{k!}
\le s^{N+1}e^{r|U|}.
\]

取期望后得到显式误差

\[
\left|\mathbb Ee^{tU}-\sum_{k=0}^{N}\frac{t^k}{k!}\mathbb EU^k\right|
\le2G(r)s^{N+1}\longrightarrow0.
\]

因此，对任意满足 \(|t|<1/q\) 的复数 \(t\)，有

\[
\mathbb E e^{tU}
=\sum_{k=0}^{\infty}\frac{t^k}{k!}\mathbb E U^k,
\qquad
\sum_{k=0}^{\infty}\frac{|t|^k}{k!}|\mathbb E U^k|
\le2G(r)<\infty.
\]

现将所得系数与 \(G(t)\) 比较。令 \(s_i\) 为 \(B\) 的奇异值。每个因子的级数由 \(c_0=1\) 及以下递推确定：

\[
c_{m+1}=\frac{2m+1}{2m+2}c_m,
\qquad
c_m=\frac{(2m)!}{4^m(m!)^2},
\qquad
(1-z)^{-1/2}=\sum_{m\ge0}c_mz^m\quad(|z|<1).
\]

这里递推来自 \((1-z)f'(z)=f(z)/2\) 的逐项系数比较；由 \(c_{m+1}/c_m<1\)，右边在 \(|z|<1\) 绝对收敛。再计算

\[
\frac{d}{dz}\bigl[(1-z)f(z)^2\bigr]
=-f(z)^2+2(1-z)f(z)f'(z)
=-f(z)^2+f(z)^2=0.
\]

在零处其值为一，所以 \((1-z)f(z)^2=1\)；因而确定了零处取值为一的分支。有限个因子相乘得到

\[
G(t)=\prod_i\left(\sum_{m\ge0}c_m s_i^{2m}t^{2m}\right),
\qquad |t|q<1.
\]

实区间内，上式和已算出的高斯积分相等。如果它的某个系数与 \(\mathbb E U^k/k!\) 不同，取最小的不同次数 \(k_0\)，两级数之差除以 \(t^{k_0}\) 后沿实数 \(t\to0\) 便得到

\[
0=\lim_{t\to0}\frac{\mathbb Ee^{tU}-G(t)}{t^{k_0}}
=\frac{\mathbb E U^{k_0}}{k_0!}-[t^{k_0}]G(t)\ne0,
\]

矛盾。因此全部系数相同；两侧在所述复圆盘内都由绝对收敛的级数给出，等式在那里也成立。指数的 \(1/k!\) 与带边标号的指数生成函数的 \(1/k!\) 一致；Wick 配对对应两侧所有块均为二元素的分划。这里仅在固定谱隙圆盘内交换收敛级数；对核心超额 \(j\) 的无穷和仍只使用形式系数恒等式。

对于任意留下贡献的图，所有度数至少二，故

\[
2(h-v)=\sum_{u}(\deg u-2)\ge0.
\]

若 \(h-v=0\)，所有顶点度数恰为二，每个连通分量都是纯二度分量。删除这些分量后，非空核心必有 \(j=h-v\ge1\)。将边标号集合分成纯二度部分与核心部分的 \(\binom kh\) 因子，和 \(1/k!\) 相乘为 \(1/[h!(k-h)!]\)，因此生成函数正好相乘。

置 \(K=BB^{\mathsf T}\)。由矩阵特征值分解和 \(-\log(1-u)=\sum_{r\ge1}u^r/r\)，

\[
\log G(t)=\frac12\sum_{r=1}^{\infty}\frac{t^{2r}}r\operatorname{tr}(K^r)
=A_st^2+B_st^4+C_st^6+O(t^8),
\]

记 \(A_s=\operatorname{tr}K/2\)、\(B_s=\operatorname{tr}(K^2)/4\)、\(C_s=\operatorname{tr}(K^3)/6\)。将指数函数按幂展开：

\[
\begin{aligned}
G(t)
&=1+(A_st^2+B_st^4+C_st^6)
+\frac12(A_st^2+B_st^4)^2+\frac16(A_st^2)^3+O(t^8)\\
&=1+A_st^2+(B_s+A_s^2/2)t^4
+(C_s+A_sB_s+A_s^3/6)t^6+O(t^8).
\end{aligned}
\]

所以

\[
\begin{aligned}
g_2&=\tfrac12\operatorname{tr}K,\\
g_4&=\tfrac14\operatorname{tr}(K^2)+\tfrac18(\operatorname{tr}K)^2,\\
g_6&=\tfrac16\operatorname{tr}(K^3)
+\tfrac18\operatorname{tr}K\operatorname{tr}(K^2)
+\tfrac1{48}(\operatorname{tr}K)^3.
\end{aligned}
\]

所有奇次系数为零，因为 \(G(t)\) 仅含 \(t^2\)。标量 \(A_s,B_s,C_s\) 仅在本段指数展开中使用。

二元素分划 \(12|34,13|24,14|23\) 共三种。两侧各选一种，共九对；其中三对相同、六对不同。相同配对时矩阵标签和为 \((\sum_{ij}B_{ij}^2)^2=(\operatorname{tr}K)^2\)；不同配对时标签和为 \(\operatorname{tr}(K^2)\)。所以纯二度贡献为

\[
\frac3{4!}(\operatorname{tr}K)^2+\frac6{4!}\operatorname{tr}(K^2)=g_4.
\]

一侧是四元素块，其 Möbius 权重为 \(-6\)；另一侧选二元素配对，有三种，每种权重为 \((-1)^2=1\)。所以行色四度项为

\[
\frac{3(-6)}{24}\sum_i\sum_{j,k}B_{ij}^2B_{ik}^2
=-\frac34\sum_i\left(\sum_jB_{ij}^2\right)^2.
\]

交换行列颜色后得到列色四度项

\[
-\frac34\sum_j\sum_{i,k}B_{ij}^2B_{kj}^2
=-\frac34\sum_j\left(\sum_iB_{ij}^2\right)^2.
\]

两侧都是四元素块，权重乘积为 \((-6)^2=36\)，故最后一项为

\[
\frac{36}{24}\sum_{i,j}B_{ij}^4=\frac32\sum_{i,j}B_{ij}^4.
\]

这四类相加得到前面的 \(F_4\)；含一度顶点的分划均因中心化而无贡献。

\Needspace{12\baselineskip}
\subsection{3.3 非高斯核心的估计}\label{ux975eux9ad8ux65afux6838ux5fc3ux7684ux4f30ux8ba1}

我们先对每条二度链的内部指标求和，再用算子范数控制所得矩阵乘积。压缩这些链后，超额为 \(j\) 的图至多有 \(2j\) 个顶点。令

\[
W=\max\left\{1,CR+\frac{C^2R^2}{1-Rq}\right\},
\qquad D=710W^3,\qquad
T_1=\frac32W^2+\frac{10}{3}W^3.
\]

\textbf{引理 3.2（压缩核心活动量）。}

在上述记号下，

\begin{equation}
\|C_j\|_R\le(Dj/n)^j\qquad(j\ge1),
\tag{3.2}
\label{eq-core-activity}
\end{equation}

且首超额满足更精确的界

\[
\|C_1\|_R\le T_1/n.
\]

其中 \(\|H\|_R=\sum_{k\ge0}|[t^k]H|R^k\)。

\textbf{证明。}

压缩所有内部顶点均为二度的极大路径。每条压缩边对应一条两端度数至少为三、长度为正的链。允许自环、重边和非连通压缩图。若 \(b\) 为度数至少为三的顶点数，\(e\) 为压缩边数，则压缩不改变超额，并给出

\[
e=b+j,\qquad
2j=\sum_{i=1}^b(d_i-2),\qquad
1\le b\le2j,\qquad e\le3j.
\]

\textbf{增饰计数恒等式。} 固定超额及原始边数，对所有有序度数序列 \((d_1,\ldots,d_b)\)、顶点的两种颜色配置、带标号半边配对及与颜色相容的正链长求和。除去临时标号后的绝对 Möbius 补偿为 \(1/(b!\prod_i d_i)\)。链收缩应先对内部数值指标求和，再取绝对值。

这里的结构顶点是两个边标号分拆的块，其数值指标取自 \([n]\)。具有相同数值指标的不同块仍视为不同顶点。原始互异性约束已由 Möbius 反演处理，因此以下直接对分拆对计数。临时为 \(b\) 个高度数顶点及每个顶点的 \(d_i\) 个半边分别加标号，并用这些标号为每条链选择唯一读取方向。遍历全部有序度数序列后，每个原始带边标号的核心有 \(b!\prod_i d_i!\) 种增饰。对一个固定的度数序列，只有与之相容的顶点标号出现；对所有序列求和才恰好产生完整的 \(b!\) 因子。每个结构顶点都能由其颜色及入射原始边标号辨认，每个入射半边则由其原始边标号辨认，故没有额外的自同构稳定子。反过来，给定增饰后的半边配对及全部正链长，所有链位置均被确定，原始 \(h\) 个边标号可以按 \(h!\) 种方式分配到这些位置。在满足二部颜色奇偶条件时，内部顶点及其颜色也随之确定。

除去原始指数生成函数的 \(h!\) 因子及临时增饰，再乘回高度数顶点的绝对 Möbius 权重 \(\prod_i(d_i-1)!\)，补偿为

\[
\frac1{b!\prod_i d_i}.
\]

这一步包含自环和链反转：将每一对不同的带标号半边排序，从较小者开始读取该链。链的反转只改变这些位置上的原始边标号次序，不另产生因子二；自环的两个半边也使用同一约定。内部二度顶点的绝对 Möbius 权重为一，也没有额外的链长阶乘。异色端点之间的链长为奇数，同色端点之间为偶数；尤其自环的长度至少为二且为偶数。计数时仍须满足这些颜色与链长条件。

例如，取一个行色四度顶点，以及两个长度为二的自环，各经过一个单独的列色二度顶点。此时 \(j=b=1\)、\(h=4\)。四个原始边标号分成两对共有三种方式，四度顶点的 Möbius 权重为 \((-1)^3 3!=-6\)，而两个二度顶点的权重相乘为 \(+1\)；除以 \(4!\) 后，其带符号贡献为

\[
-\frac34\sum_i\left(\sum_k B_{ik}^2\right)^2.
\]

即使某些求和项中数值指标相等，两个结构列顶点仍然不同。交换两种颜色得到 \(-\tfrac34\sum_k(\sum_i B_{ik}^2)^2\)。绝对补偿界使用的是这些项的绝对值。重边的例子是一个行顶点和一个列顶点之间的三条长度为一的链，其贡献为 \(\tfrac{(2!)^2}{3!}\sum_{i,k}B_{ik}^3=\tfrac23\sum_{i,k}B_{ik}^3\)，恰为前述三次项。上述自环和平行边的计数均遵守反转约定，压缩前的图仍为二部图。

仅在取上界时忽略颜色限制。固定 \(j,b\) 后，所得总补偿为

\[
U_{j,b}=
\frac{2^b}{b!}\frac{(2e)!}{2^e e!}
\sum_{\substack{d_1+\cdots+d_b=2e\\d_i\ge3}}
\frac1{d_1\cdots d_b},
\qquad e=b+j.
\]

以下估计须先对每条链的内部数值指标求和，再取绝对值；若先逐项将 \(B\) 换成 \(|B|\)，便无法保留算子谱隙的控制。长度为一的链对应矩阵元素，绝对值不超过 \(C/n\)。对于长度 \(\ell\ge2\) 的链，先求和内部标签，其收缩是由 \(B,B^{\mathsf T}\) 交替相乘所得矩阵的一个元素。两端行或列的二范数不超过 \(C/\sqrt n\)，中间因子的算子范数不超过 \(q\)。因此，收缩的绝对值不超过

\[
\frac{C^2}{n}q^{\ell-2}.
\]

按 \(R\) 加权对长度求和，每条链的总活动量不超过 \(W/n\)。对 \(b\) 个高度数顶点标签求和至多给出 \(n^b\)，故

\[
\|C_j\|_R\le n^{-j}
\sum_{b=1}^{2j}U_{j,b}W^{b+j}.
\]

度数组合数为 \(\binom{2j-1}{b-1}\le4^j\)：写 \(d_i=3+x_i\)，则 \(\sum_i x_i=2j-b\)。此外，\(\prod_i d_i^{-1}\le3^{-b}\)，半边配对数不超过 \((6j)^{b+j}\)。于是

\[
\sum_{b=1}^{2j}U_{j,b}W^{b+j}
\le(24jW)^j
\sum_{b=1}^{2j}\frac{(4jW)^b}{b!}.
\]

取 \(\theta=1/(2W)\le1/2\)，则

\[
\sum_{b=0}^{2j}\frac{(4jW)^b}{b!}
\le\theta^{-2j}\exp(4jW\theta)
=(2\mathrm e W)^{2j}.
\]

因此，活动量至多为 \((96\mathrm e^2W^3j)^j\le(710W^3j)^j\)。后一不等式使用 \(\mathrm e<87/32\)；将指数级数求和至六次，再以 \(8/(7\cdot7!)\) 控制尾项，即得此界。

当 \(j=1\) 时仅有 \(b=1,2\)。补偿公式分别给出

\[
U_{1,1}=\frac32,\qquad U_{1,2}=\frac{10}{3}.
\]

结合链界即有 \(\|C_1\|_R\le T_1/n\)。

压去一个二度顶点同时减少一条边和一个顶点，故 \(h-v\) 不变。压缩图满足

\[
2e=\sum_{i=1}^bd_i,\quad e-b=j,\quad
2j=\sum_{i=1}^b(d_i-2)\ge b,
\]

因此 \(b\le2j\)、\(e=b+j\le3j\)。除去临时标号并计入 Möbius 权重后，所得系数为

\[
\frac1{h!}\cdot\frac{h!}{b!\prod_i d_i!}\cdot\prod_i(d_i-1)!
=\frac1{b!}\prod_i\frac{(d_i-1)!}{d_i!}
=\frac1{b!\prod_i d_i}.
\]

对 \(2e\) 个半边，配对数由递推 \(p_{2e}=(2e-1)p_{2e-2}\)、\(p_0=1\) 得

\[
p_{2e}=(2e-1)(2e-3)\cdots1=\frac{(2e)!}{2^ee!}.
\]

颜色配置数为 \(2^b\)；故在忽略不相容颜色和链长后得到 \(U_{j,b}\)。链长一的收缩为 \(B_{ij}\)。若长度 \(\ell\ge2\)，写收缩为 \(r^{\mathsf T}M_1\cdots M_{\ell-2}c\)，其中 \(r,c\) 是 \(B\) 的行或列，则

\[
\|r\|_2^2,\|c\|_2^2\le n(C/n)^2=C^2/n,
\]

\[
|r^{\mathsf T}M_1\cdots M_{\ell-2}c|
\le\|r\|_2\prod_a\|M_a\|_{\rm op}\|c\|_2
\le(C^2/n)q^{\ell-2}.
\]

按长度加权后

\[
R\frac Cn+\sum_{\ell=2}^{\infty}R^\ell\frac{C^2}nq^{\ell-2}
=\frac1n\left(CR+\frac{C^2R^2}{1-Rq}\right)\le W/n.
\]

剩余 \(b\) 个顶点的数值标签至多有 \(n^b\) 种选择，故 \(n^b(W/n)^e=n^{-j}W^{b+j}\)。链的内部求和保留了因子 \(q^{\ell-2}\)。

度数组合计数：置 \(x_i=d_i-3\ge0\)，则 \(\sum_ix_i=2j-b\)，所以

\[
\#\{(d_1,\ldots,d_b)\}=\binom{(2j-b)+b-1}{b-1}
=\binom{2j-1}{b-1}\le2^{2j-1}\le4^j.
\]

并且 \(\prod_i1/d_i\le3^{-b}\)，配对数至多为 \((2e)^e\le(6j)^{b+j}\)。于是

\[
\begin{aligned}
U_{j,b}W^{b+j}
&\le\frac{2^b}{b!}(6j)^{b+j}4^j3^{-b}W^{b+j}\\
&=(24jW)^j\frac{(4jW)^b}{b!}.
\end{aligned}
\]

若 \(0<\theta\le1\) 且 \(b\le2j\)，则 \(\theta^b\ge\theta^{2j}\)，故

\[
\sum_{b=0}^{2j}\frac{(4jW)^b}{b!}
\le\theta^{-2j}\sum_{b=0}^{\infty}\frac{(4jW\theta)^b}{b!}
=\theta^{-2j}e^{4jW\theta}.
\]

取 \(\theta=1/(2W)\) 后，上式为 \((2W)^{2j}e^{2j}\)。乘以 \((24jW)^j\)，即得 \((96e^2W^3j)^j\)。这里用到的有理数界为

\[
e<\sum_{h=0}^6\frac1{h!}+\frac1{7!}\sum_{r=0}^{\infty}8^{-r}
=\frac{1957}{720}+\frac1{4410}<\frac{87}{32},
\]

\[
96(87/32)^2=\frac{22707}{32}=709.59375<710.
\]

当 \(j=1\)：\(b=1,e=2,d_1=4\)，所以

\[
U_{1,1}=\frac2{1!}\frac{4!}{2^2\,2!}\frac14=\frac32.
\]

当 \(b=2,e=3\)，唯一度数列为 \((3,3)\)，故

\[
U_{1,2}=\frac{2^2}{2!}\frac{6!}{2^3\,3!}\frac1{3\cdot3}=\frac{10}{3}.
\]

合并两种情形，得 \(T_1=(3/2)W^2+(10/3)W^3\)。

固定一个压缩核心，给剩下的 \(b\) 个高度数顶点指定数值标签。设其 \(e\) 条压缩边对应的内部结构顶点集合分别为 \(V_1,\ldots,V_e\)；这些集合互不相交。数值标签可以相等，但每个结构顶点对应独立求和变量。因此有限和的分配律给出

\[
\sum_{\text{全部内部标签}}\prod_{r=1}^e
\left(\prod_{a\in\mathrm{chain}_r}B_{i_aj_a}\right)
=\prod_{r=1}^e\left[
\sum_{\text{链 }r\text{ 的内部标签}}
\prod_{a\in\mathrm{chain}_r}B_{i_aj_a}\right].
\]

每个中括号是一个交替矩阵乘积的单个元素；现在才对它取绝对值。正链长的求和也分离为 \(e\) 个几何级数。因此固定高度数标签后的总绝对活动量至多为 \((W/n)^e\)。再对 \(n^b\) 种高度数标签求和，得到

\[
n^b(W/n)^e=W^en^{b-e}=W^{b+j}n^{-j}.
\]

这里各链的内部结构顶点集互不相交，故求和变量独立；数值标签仍可相同。

当 \(j=1,b=1\) 时，有一个四度顶点、四个半边。配对只能为

\[
(12)(34),\quad(13)(24),\quad(14)(23).
\]

每个配对给两个自环，颜色有两种，补偿分母 \(b!d_1=4\)；总数 \(2\cdot3/4=3/2\)。

当 \(j=1,b=2\) 时，两顶点各三度。半边配对分成两类：全部三条边连接两个顶点，数目为 \(3!=6\)；每个顶点各有一个自环，剩下的两个半边连成一条边，数目为 \(3\cdot3=9\)。两类和为十五。顶点颜色四种，补偿分母为 \(2!\cdot3\cdot3=18\)，所以

\[
\frac4{18}(6+9)=\frac{10}{3}.
\]

忽略颜色与链长的相容条件给出上界，两个情形分别产生 \(T_1\) 的两项。 \(\square\)

\Needspace{12\baselineskip}
\subsection{3.4 线性系数窗口与阶乘恢复}\label{ux7ebfux6027ux7cfbux6570ux7a97ux53e3ux4e0eux9636ux4e58ux6062ux590d}

活动量界\hyperref[eq-core-activity]{式 (3.2)}在 \(j/n\) 较小时有效。取

\[
\alpha=\min\left\{\frac14,\frac1{16D},\frac{\log\sigma}{2}\right\},
\qquad M=\lfloor\alpha n\rfloor.
\]

并令 \(v_j=(Dj/n)^j\)。当 \(j+1\le M\) 时，

\[
\frac{v_{j+1}}{v_j}
=\frac{D(j+1)}n\left(1+\frac1j\right)^j
\le\frac{\mathrm e D(j+1)}n
\le3D\alpha<\frac12.
\]

故对 \(n\ge4/\alpha\)，

\[
\sum_{j=2}^{M}v_j\le2v_2=\frac{8D^2}{n^2}.
\]

形式核心恒等式中，次数不超过 \(M\) 的贡献只能来自超额 \(j\le M\)。结合 \(G\) 系数非负、活动量引理及加权系数和的卷积不等式，得到

\[
\sum_{k=0}^{M}|F_k-[t^k]G|R^k
\le\mathcal G_{C,q}(R)\left(\frac{T_1}n+\frac{8D^2}{n^2}\right).
\]

该估计仅涉及次数不超过 \(M\) 的系数。

记 \(g_k=[t^k]G\)，并在 \(0\le k\le M\) 时令 \(r_k=n^k/(n)_k\)。由于 \(\alpha\le1/4\)，

\[
\log r_k
=\sum_{\ell=0}^{k-1}-\log(1-\ell/n)
\le\frac{k(k-1)}{2n(1-\alpha)}
\le\frac{k^2}n
\le\alpha k.
\]

由 \(\alpha\le(\log\sigma)/2\) 可知 \(r_k\le\sigma^k\)，并且

\[
0\le r_k-1
\le\frac{k(k-1)}{2n(1-\alpha)}r_k.
\]

因为 \(a_k=r_kF_k\)，可以分解为

\[
a_k-g_k=r_k(F_k-g_k)+(r_k-1)g_k.
\]

第一部分由上述 \(R\) 加权估计控制。对第二部分，使用高斯系数的二阶矩：

\[
\sum_{k\ge0}k(k-1)g_k\sigma^k
=\left((t\partial_t)^2-t\partial_t\right)G(t)\big|_{t=\sigma}.
\]

令 \(a_i=\sigma^2s_i^2\)。对 \(G(t)=\prod_i(1-t^2s_i^2)^{-1/2}\) 求导，得到

\[
\frac{\left((t\partial_t)^2-t\partial_t\right)G(t)}{G(t)}
\bigg|_{t=\sigma}
=\left(\sum_i\frac{a_i}{1-a_i}\right)^2
+\sum_i\frac{a_i(1+a_i)}{(1-a_i)^2}.
\]

定义

\[
u_\sigma=\frac{\sigma^2C^2}{1-\sigma^2q^2},
\qquad
v_\sigma=\frac{\sigma^2C^2(1+\sigma^2q^2)}
{(1-\sigma^2q^2)^2},
\]

\[
K_G=\frac{\mathcal G_{C,q}(\sigma)}{2(1-\alpha)}
\left(u_\sigma^2+v_\sigma\right).
\]

导数公式中的两项分别不超过 \(u_\sigma^2\) 与 \(v_\sigma\)。因此，

\[
\sum_{k=0}^{M}(r_k-1)g_k\le K_G/n.
\]

这给出误差界的前两项。

由 \(e<3\)、\(\alpha\le1/(16D)\)，

\[
\frac{v_{j+1}}{v_j}\le3D\alpha\le3/16<1/2,
\quad
\sum_{j=2}^Mv_j\le v_2\sum_{r=0}^{\infty}2^{-r}
=2(2D/n)^2=8D^2/n^2.
\]

若 \(H(t)=\sum h_kt^k\)、\(K(t)=\sum k_kt^k\)，则

\[
\sum_{r=0}^M|[t^r](HK)|R^r
\le\sum_{a,b\ge0}|h_a||k_b|R^{a+b}
=\|H\|_R\|K\|_R.
\]

在 \(r\le M\) 的系数中只需 \(j\le M\)；所以先截取有限超额和再使用这个不等式，没有计算 \(\sum_{j\ge1}C_j(R)\)。

阶乘恢复中的积分界为

\[
-\log(1-u)=\int_0^u\frac{dv}{1-v}\le\frac{u}{1-\alpha}
\quad(0\le u\le\alpha).
\]

累加 \(u=0,1/n,\ldots,(k-1)/n\) 得

\[
\log r_k\le\frac{k(k-1)}{2n(1-\alpha)}\le k^2/n\le\alpha k,
\quad r_k\le e^{\alpha k}\le\sigma^k.
\]

由 \(e^x-1\le xe^x\) 得 \(r_k-1\le[k(k-1)/(2n(1-\alpha))]r_k\)。因此

\[
\begin{aligned}
\sum_{k=0}^M|a_k-g_k|
&\le\sum_{k=0}^Mr_k|F_k-g_k|+\sum_{k=0}^M(r_k-1)g_k\\
&\le\sum_{k=0}^MR^k|F_k-g_k|
+\frac1{2n(1-\alpha)}\sum_{k\ge0}k(k-1)g_k\sigma^k.
\end{aligned}
\]

令 \(D_t=t\partial_t\)、\(a_i(t)=t^2s_i^2\)。逐次微分为

\[
\begin{aligned}
D_t\log G&=\sum_i\frac{a_i}{1-a_i}=:A(t),\\
D_tA&=\sum_i\frac{2a_i}{(1-a_i)^2},\\
D_tG&=GA,\\
D_t^2G&=G\left[A^2+\sum_i\frac{2a_i}{(1-a_i)^2}\right],\\
\frac{(D_t^2-D_t)G}{G}
&=A^2+\sum_i\left[\frac{2a_i}{(1-a_i)^2}-\frac{a_i}{1-a_i}\right]\\
&=A^2+\sum_i\frac{a_i(1+a_i)}{(1-a_i)^2}.
\end{aligned}
\]

在 \(t=\sigma\)，\(a_i\le\sigma^2q^2\)，故

\[
A(\sigma)\le\frac{\sigma^2C^2}{1-\sigma^2q^2}=u_\sigma,
\quad
\sum_i\frac{a_i(1+a_i)}{(1-a_i)^2}
\le\frac{\sigma^2C^2(1+\sigma^2q^2)}{(1-\sigma^2q^2)^2}=v_\sigma.
\]

结合 \(G(\sigma)\le\mathcal G_{C,q}(\sigma)\)，恰得到 \(K_G/n\)。

定义投影 \(\mathcal P_M\sum_kh_kt^k=\sum_{k=0}^Mh_kt^k\)。因为超额 \(j\) 的非空核心至少有 \(j+1\) 条原始边，次数不超过 \(M\) 时不会出现 \(j>M\)，故有

\[
\mathcal P_M(F-G)=\mathcal P_M\left(G\sum_{j=1}^MC_j\right).
\]

这右边是有限个核心级数之和，因而加权范数合法。由三角不等式和卷积估计，

\[
\begin{aligned}
\sum_{k=0}^M|F_k-g_k|R^k
&=\|\mathcal P_M(F-G)\|_R\\
&\le\left\|G\sum_{j=1}^MC_j\right\|_R\\
&\le\|G\|_R\sum_{j=1}^M\|C_j\|_R\\
&\le\mathcal G_{C,q}(R)\left[T_1/n+\sum_{j=2}^M(Dj/n)^j\right]\\
&\le\mathcal G_{C,q}(R)(T_1/n+8D^2/n^2).
\end{aligned}
\]

此处只用有限个 \(C_j\) 的加权范数。恢复 \(a_k\) 时，第一部分满足

\[
\sum_{k=0}^Mr_k|F_k-g_k|
\le\sum_{k=0}^M\sigma^k|F_k-g_k|
\le\sum_{k=0}^MR^k|F_k-g_k|.
\]

第二部分 \((r_k-1)g_k\) 由前述高斯导数矩估计。

\Needspace{12\baselineskip}
\subsection{3.5 控制剩余系数}\label{ux63a7ux5236ux5269ux4f59ux7cfbux6570}

前一节已估计次数不超过 \(M\) 的系数，且 \(M\) 与 \(n\) 成正比。为用 Cauchy 不等式估计更高次系数，只需在固定圆周 \(|t|=R>1\) 上将 \(f\) 控制为 \(\exp(O(\sqrt n))\)。以下两个引理给出这一上界，并保留归一化因子 \(n!/n^n\)。

\textbf{引理 3.3（永久式极化）。}

对复方阵 \(Z\)，其奇异值取 \(Z^*Z\) 的特征值的非负平方根，相应范数遵循前面的约定。记 \(|Z|=(Z^*Z)^{1/2}\) 为算子绝对值，其中平方根取 Hermitian 半正定平方根，而非逐元素取绝对值。对 Hermitian 矩阵写 \(H\ge0\)，意思是每个向量均满足 \(x^*Hx\ge0\)。则

\[
|\operatorname{per}Z|
\le\sqrt{\operatorname{per}|Z|\,
\operatorname{per}|Z^*|}.
\]

\textbf{证明。}

在复向量空间采用内积 \(\langle x,y\rangle=x^*y\)。在 \((\mathbb C^n)^{\otimes n}\) 中令

\[
\xi=(n!)^{-1/2}
\sum_{\pi\in \mathfrak S_n}
e_{\pi(1)}\otimes\cdots\otimes e_{\pi(n)}.
\]

则 \(\|\xi\|=1\)，且 \(\langle\xi,Z^{\otimes n}\xi\rangle=\operatorname{per}Z\)。 作极分解 \(Z=U|Z|\)；若 \(Z\) 奇异，将极分解的部分等距因子延拓为酉矩阵 \(U\)。记 \(T=|Z|^{\otimes n}\)，\(V=U^{\otimes n}\)。在由半正定算子 \(T\) 诱导的半内积中应用 Cauchy--Schwarz 不等式，得到

\[
|\langle\xi,VT\xi\rangle|
\le\sqrt{\langle V^*\xi,TV^*\xi\rangle
\langle\xi,T\xi\rangle}.
\]

两个因子分别为 \(\operatorname{per}(U|Z|U^*)=\operatorname{per}|Z^*|\) 与 \(\operatorname{per}|Z|\)。由于它们是半正定张量幂的对角矩阵元，均非负。\(\square\)

\textbf{引理 3.4（半正定永久式界）。}

设 \(n\ge1\)，\(H\) 为 \(n\times n\) 的 Hermitian 半正定矩阵，特征值为 \(\lambda_1,\ldots,\lambda_n\)。则

\[
\operatorname{per}H\le\frac{n!}{n^n}h_n(\lambda_1,\ldots,\lambda_n),
\]

其中 \(h_n\) 为 \(n\) 次完全齐次对称多项式，定义为

\[
h_n(\lambda_1,\ldots,\lambda_n)
=\sum_{\substack{k_1,\ldots,k_n\ge0\\k_1+\cdots+k_n=n}}
\prod_{i=1}^n\lambda_i^{k_i}.
\]

各 \(k_i\) 为整数；多项式中的零次幂恒为一，包括 \(\lambda_i=0\) 时的 \(\lambda_i^0\)。这里的记号只用于本引理。若 \(\lambda_1=1\) 且 \(0\le\lambda_i<1\)（\(i\ge2\)），则

\[
\operatorname{per}H\le\frac{n!}{n^n}
\prod_{i=2}^n(1-\lambda_i)^{-1}.
\]

\textbf{证明。}

令 \(Z\) 为协方差矩阵为 \(H\) 的圆对称复高斯向量。复 Wick 公式给出 \(\operatorname{per}H=\mathbb E\prod_i|Z_i|^2\)。 逐点 AM--GM 不等式把该乘积控制为 \(n^{-n}\|Z\|^{2n}\)。酉对角化后，\(\|Z\|^2\) 与 \(\sum_i\lambda_i|\zeta_i|^2\) 同分布，其中 \(\zeta_i\) 为独立标准圆对称复高斯变量，其矩为 \(\mathbb E|\zeta_i|^{2k}=k!\)。于是，多项式展开给出

\[
\mathbb E\left(\sum_i\lambda_i|\zeta_i|^2\right)^n
=n!h_n(\lambda_1,\ldots,\lambda_n).
\]

最后，\(h_n(1,\lambda_2,\ldots,\lambda_n)\) 是关于 \(\lambda_2,\ldots,\lambda_n\)、总次数不超过 \(n\) 的全部单项式之和，不超过各自无穷几何级数的乘积。\(\square\)

引理 3.4 中的高斯恒等式和 AM--GM 比较，是 Han--Niles-Weed {[}3，arXiv v2，引理 4.3--4.4{]} 用于半正定输入的工具。引理 3.3 将它们连接到这里需要的一般矩阵。

令 \(\beta=RC/(1-Rq)\)，并记 \(P_n=J_n/n\)，它是到全一向量方向的正交投影。一般写 \(u^\perp=\{v:u^*v=0\}\) 表示向量 \(u\) 的正交补。双中心化假设给出 \(P_nB=BP_n=0\)。在 \(|t|=R\) 上，因而有

\[
|P_n+tB|=P_n+R(B^{\mathsf T}B)^{1/2},
\qquad
|(P_n+tB)^*|=P_n+R(BB^{\mathsf T})^{1/2}.
\]

两者的特征值均为全一方向上的 \(1\)，以及 \(\mathbf1^\perp\) 上的 \(Rs_i\)，必要时包含零奇异值。由于 \(Rq<1\)，两个引理结合精确恒等式 \(f(t)=(n^n/n!)\operatorname{per}(P_n+tB)\)，给出

\[
|f(t)|\le\prod_i(1-Rs_i)^{-1}
\le\exp\!\left(\frac{R}{1-Rq}\sum_i s_i\right)
\le\exp(\beta\sqrt n).
\]

最后一步使用 \(\sum_i s_i\le\sqrt n\,\|B\|_{\mathrm F}\le C\sqrt n\)。因此，Cauchy 系数估计给出

\[
\sum_{k>M}|a_k|
\le\frac{R}{R-1}
\exp\!\left(\beta\sqrt n-\alpha n\log R\right).
\]

对于高斯尾部，利用其系数非负，可用稍松但便利的界

\[
\sum_{k>M}g_k
\le \mathcal G_{C,q}(\sigma)\sigma^{-M}
\le\sigma\mathcal G_{C,q}(\sigma)\exp\!\left(-\alpha n\log\sigma\right).
\]

两个尾项均随 \(n\) 指数衰减。将它们与有限窗口中的 \(O(n^{-1})\) 估计合并，即证明定理 3.1。

对任意 \(M\)，

\[
\begin{aligned}
\langle\xi,M^{\otimes n}\xi\rangle
&=\frac1{n!}\sum_{\pi,\sigma\in \mathfrak S_n}\prod_{h=1}^nM_{\pi(h),\sigma(h)}\\
&=\frac1{n!}\sum_{\pi\in \mathfrak S_n}\sum_{\tau\in \mathfrak S_n}\prod_{i=1}^nM_{i,\tau(i)}
=\operatorname{per}M.
\end{aligned}
\]

这里 \(\tau=\sigma\circ\pi^{-1}\)；固定 \(\pi\) 时 \(\sigma\mapsto\tau\) 是双射。\(n!\) 个基张量相互正交，故 \(\|\xi\|^2=(n!)^{-1}n!=1\)。若 \(T\ge0\)，

\[
|\langle a,Tb\rangle|^2
=|\langle T^{1/2}a,T^{1/2}b\rangle|^2
\le\langle a,Ta\rangle\langle b,Tb\rangle.
\]

取 \(a=V^*\xi,b=\xi\)，利用 \(VTV^*=(U|Z|U^*)^{\otimes n}=|Z^*|^{\otimes n}\)，得到引理 3.3。

标准圆对称复变量 \(\zeta=x+iy\) 的密度为 \(\pi^{-1}e^{-|z|^2}\)。极坐标积分给出

\[
\begin{aligned}
\mathbb E|\zeta|^{2k}
&=\pi^{-1}\int_0^{2\pi}\!d\theta\int_0^{\infty}r^{2k}e^{-r^2}r\,dr\\
&=\int_0^{\infty}u^ke^{-u}\,du=k!.
\end{aligned}
\]

记 \(I_k=\int_0^\infty u^ke^{-u}\,du\)，分部积分给出 \(I_k=kI_{k-1}\)、\(I_0=1\)。

对整数 \(a,b\ge0\)，极坐标积分给出一维混合矩

\[
\begin{aligned}
\mathbb E[\zeta^a(\overline\zeta)^b]
&=\frac1\pi\int_0^\infty r^{a+b+1}e^{-r^2}\,dr
\int_0^{2\pi}e^{i(a-b)\theta}\,d\theta\\
&=\begin{cases}a!,&a=b,\\0,&a\ne b.\end{cases}
\end{aligned}
\]

若 \(a\ne b\)，角积分为 \([e^{i(a-b)\theta}/(i(a-b))]_0^{2\pi}=0\)；若 \(a=b\)，角积分为 \(2\pi\)，再用 \(u=r^2\) 得 \(\int_0^\infty u^ae^{-u}\,du=a!\)。

令 \(\zeta_1,\ldots,\zeta_n\) 独立，\(Z=L\zeta\)、\(H=LL^*\)。对指标列 \(\boldsymbol a,\boldsymbol b\in[n]^n\)，上式和独立性表明：当某个数值指标在两列中的重数不同时，混合矩为零；重数均相同且为 \(c_j\) 时，混合矩为 \(\prod_jc_j!\)。这正好等于把第一列的每个位置与第二列的同值位置双射配对的方法数。因此

\[
\mathbb E\prod_{i=1}^n\zeta_{a_i}\prod_{j=1}^n\overline\zeta_{b_j}
=\sum_{\sigma\in \mathfrak S_n}\prod_{i=1}^n\delta_{a_i,b_{\sigma(i)}}.
\]

展开乘积并交换有限求和，得到

\[
\begin{aligned}
\mathbb E\prod_{i=1}^n|Z_i|^2
&=\sum_{\boldsymbol a,\boldsymbol b}
\prod_iL_{i,a_i}\prod_j\overline L_{j,b_j}
\sum_{\sigma\in \mathfrak S_n}\prod_i\delta_{a_i,b_{\sigma(i)}}\\
&=\sum_{\sigma\in \mathfrak S_n}\sum_{\boldsymbol a}
\prod_i L_{i,a_i}\overline L_{\sigma(i),a_i}\\
&=\sum_{\sigma\in \mathfrak S_n}\prod_i
\left(\sum_{a=1}^nL_{i,a}\overline L_{\sigma(i),a}\right)\\
&=\sum_{\sigma\in \mathfrak S_n}\prod_iH_{i,\sigma(i)}
=\operatorname{per}H.
\end{aligned}
\]

第二行使用 \(b_{\sigma(i)}=a_i\) 消去全部 \(b\) 指标，第三行使用各个 \(a_i\) 独立求和。\(L\) 可以奇异；取 \(L=H^{1/2}\) 即可，因此本证明涵盖半正定矩阵有零特征值的情形。协方差的含义也由同样的二次矩计算确定：\(\mathbb EZ_i\overline Z_j=H_{ij}\)，\(\mathbb EZ_iZ_j=0\)。

对每个样本，\(|Z_i|^2\ge0\)；AM--GM 给出

\[
\prod_i|Z_i|^2\le\left(\frac{\sum_i|Z_i|^2}{n}\right)^n,
\qquad \operatorname{per}H\le n^{-n}\mathbb E\|Z\|_2^{2n}.
\]

酉变换保持 \(\|Z\|\)，再由多项式定理，

\[
\begin{aligned}
\mathbb E\left(\sum_i\lambda_i|\zeta_i|^2\right)^n
&=\sum_{k_1+\cdots+k_n=n}\frac{n!}{\prod_i k_i!}
\prod_i\lambda_i^{k_i}\mathbb E|\zeta_i|^{2k_i}\\
&=n!\sum_{k_1+\cdots+k_n=n}\prod_i\lambda_i^{k_i}
=n!h_n(\lambda).
\end{aligned}
\]

若 \(\lambda_1=1\)，以 \(k_1=n-\sum_{i\ge2}k_i\) 消去第一变量，

\[
h_n(1,\lambda_2,\ldots,\lambda_n)
=\sum_{\substack{k_2,\ldots,k_n\ge0\\\sum k_i\le n}}\prod_{i\ge2}\lambda_i^{k_i}
\le\prod_{i=2}^n\sum_{k_i=0}^{\infty}\lambda_i^{k_i}
=\prod_{i=2}^n(1-\lambda_i)^{-1}.
\]

记 \(P=P_n\)。因 \(P^2=P\)、\(PB=B^{\mathsf T}P=BP=PB^{\mathsf T}=0\)，

\[
(P+tB)^*(P+tB)=P+R^2B^{\mathsf T}B
\quad(|t|=R).
\]

这两个部分作用在正交子空间上，故平方根为 \(P+R(B^{\mathsf T}B)^{1/2}\)。两个极化因子具有相同特征值 \(1,Rs_1,\ldots\)；全一方向对应的零奇异值可从乘积中省去，因为相应因子为一。定义 \(A_R=\prod_i(1-Rs_i)^{-1}>0\)。只对这个正上界取对数，避免 \(f(t)=0\) 时出现未定义的 \(\log|f(t)|\)。代入归一化因子，得

\[
\begin{aligned}
|f(t)|&=\frac{n^n}{n!}|\operatorname{per}(P+tB)|
\le\prod_i(1-Rs_i)^{-1},\\
\log A_R&=-\sum_i\log(1-Rs_i)\le\sum_i\frac{Rs_i}{1-Rs_i}
\le\frac R{1-Rq}\sum_i s_i
\le\frac{RC}{1-Rq}\sqrt n=\beta\sqrt n.
\end{aligned}
\]

所以 \(|f(t)|\le A_R\le e^{\beta\sqrt n}\)，包括 \(f(t)=0\) 的情形。零点并非纯粹假想：取 \(n=3\)、

\[
E=\begin{pmatrix}0&1&-1\\-1&0&1\\1&-1&0\end{pmatrix},\qquad B=E/3.
\]

此时双侧中心化，且

\[
B^{\mathsf T}B=\frac19(3I_3-J_3).
\]

在全一方向上 \(J_3\) 的特征值为三，在零和子空间上为零，故 \(B^{\mathsf T}B\) 的特征值为 \(0,1/3,1/3\)。因此 \(\|B\|_{\rm op}=1/\sqrt3\le3/5=:q<1\)。六个永久式项分别来自恒等置换、三个换位和两个三循环，所以

\[
\begin{aligned}
6f(t)&=1+3(1-t^2)+(1+t)^3+(1-t)^3\\
&=1+3-3t^2+2+6t^2=6+3t^2,\\
f(i\sqrt2)&=1+(i\sqrt2)^2/2=0.
\end{aligned}
\]

取 \(R=\sqrt2<5/3=1/q\) 和 \(\sigma=(1+\sqrt2)/2\)，则 \(f\) 在所取圆周 \(|t|=R\) 上有零点，故 \(\log|f(t)|\) 在该圆周上并非处处有定义。

Cauchy 系数公式沿圆周 \(t=Re^{i\theta}\) 写成

\[
a_k=\frac1{2\pi R^k}\int_0^{2\pi}f(Re^{i\theta})e^{-ik\theta}\,d\theta,
\qquad |a_k|\le e^{\beta\sqrt n}R^{-k}.
\]

于是

\[
\sum_{k>M}|a_k|\le e^{\beta\sqrt n}\frac{R^{-M}}{R-1}
\le\frac R{R-1}e^{\beta\sqrt n-\alpha n\log R}.
\]

这里 \(M\ge\alpha n-1\)。另一侧系数非负，所以

\[
\sum_{k>M}g_k\le\sigma^{-M}\sum_{k>M}g_k\sigma^k
\le\sigma\mathcal G_{C,q}(\sigma)e^{-\alpha n\log\sigma}.
\]

令 \(c=\alpha\log R>0\)。若 \(n\ge(2\beta/c)^2\)，则 \(\beta\sqrt n-cn\le-cn/2\)。由 \(ne^{-cn/2}\to0\)，第一个尾项为 \(o(n^{-1})\)，第二个同理。此处用到 \(M=\lfloor\alpha n\rfloor\) 随 \(n\) 线性增长。

\Needspace{12\baselineskip}
\subsection{3.6 合并误差界}\label{ux5408ux5e76ux8befux5deeux754c}

\textbf{命题 3.5（显式误差界）。}

对每个整数 \(n\ge4/\alpha\)，

\[
\begin{aligned}
|f(1)-G(1)|
&\le \frac{\mathcal G_{C,q}(R)T_1+K_G}{n}
+\frac{8\mathcal G_{C,q}(R)D^2}{n^2}\\
&\quad+\frac{R}{R-1}\exp\!\left(\beta\sqrt n-\alpha n\log R\right)\\
&\quad+\sigma\mathcal G_{C,q}(\sigma)\exp\!\left(-\alpha n\log\sigma\right).
\end{aligned}
\]

前两项来自第 3.4 节的系数窗口及阶乘恢复误差，后两项来自第 3.5 节的尾界。两个指数项均为 \(o(n^{-1})\)，从而得到定理 3.1 的一致误差率。

令 \(c_1=\alpha\log R>0\)、\(c_2=\alpha\log\sigma>0\)、\(\beta=RC/(1-Rq)\ge0\)。对

\[
n\ge N_1:=\max\{1,(2\beta/c_1)^2\},
\]

有

\[
\sqrt n\ge2\beta/c_1
\ \Longrightarrow\ {}
\beta\sqrt n\le(c_1/2)n
\ \Longrightarrow\ {}
\beta\sqrt n-c_1n\le-c_1n/2.
\]

对任意固定 \(c>0\)，函数 \(u(x)=xe^{-cx}\) 在 \(x\ge0\) 上满足

\[
u'(x)=e^{-cx}(1-cx),\qquad
\max_{x\ge0}u(x)=u(1/c)=1/(ec).
\]

因此，对 \(n\ge N_1\)，

\[
\frac R{R-1}\exp(\beta\sqrt n-c_1n)
\le\frac R{R-1}e^{-c_1n/2}
\le\frac{2R}{ec_1(R-1)}\frac1n,
\]

\[
\sigma\mathcal G_{C,q}(\sigma)e^{-c_2n}
\le\frac{\sigma\mathcal G_{C,q}(\sigma)}{ec_2}\frac1n.
\]

若还要求 \(n\ge4/\alpha\)，命题 3.5 的右边便至多为

\[
\frac1n\left[
\mathcal G_{C,q}(R)T_1+K_G+8\mathcal G_{C,q}(R)D^2
+\frac{2R}{ec_1(R-1)}
+\frac{\sigma\mathcal G_{C,q}(\sigma)}{ec_2}\right].
\]

括号中的每个量都只依赖先固定的 \(C,q,R,\sigma\)，没有依赖正在估计的矩阵。若要把界写成所有 \(n\ge1\) 都成立的同一个常数，还须覆盖有限个小阶数。由 \(|1+E_{ij}|\le1+C\)，

\[
|f(1)|\le\frac1{n!}\sum_{\pi\in \mathfrak S_n}(1+C)^n=(1+C)^n,
\qquad 1\le G(1)\le\mathcal G_{C,q}(1).
\]

令 \(N=\lceil\max\{N_1,4/\alpha\}\rceil\)，再把统一常数增大到至少

\[
\max_{1\le n<N}n\left[(1+C)^n+\mathcal G_{C,q}(1)\right]
\]

即可覆盖有限个小阶数；若最大值集合为空，约定其值为零。 \(\square\)

\textbf{已有永久式证明中的定量余项。}

Divoux--Wang {[}12, 命题 3.1{]} 对非负双随机矩阵给出了统一的 \(1+o(1)\) 永久式近似。对于本文计数中出现的这一类矩阵，其证明可以细化为 \(1+O(N^{-1})\)。这里的误差阶不是原命题的陈述，而是保留其估计中的参数依赖后得到的改进。本文定理 3.1 允许任意实双中心化矩阵 \(E\)，不要求 \(J+E\) 非负；我们保留这一形式及其详细证明，不将两篇论文的定理假设视为完全相同。 定理假设的表述差别并不构成优先性主张，也不声称超越 McCullagh {[}5, §4.1{]} 的适用范围；另见主文第 3 节。

上述细化只涉及 {[}12, 第 3.2--3.3 节{]} 的两处估计。其有限系数和已被控制为

\[
\sum_{1\le E\le K\sqrt N}
\frac{(cE/N)^{\lceil E/3\rceil}}{1-cE/N}.
\]

令 \(\eta_N=cK/\sqrt N\)。当 \(N\) 充分大时，这个和至多为

\[
\frac{c}{N(1-\eta_N)}
\sum_{E\ge1}E\eta_N^{\lceil E/3\rceil-1}
=O(N^{-1}).
\]

截断之后的系数尾项仍由原文的指数估计控制。在其 Gamma 平均步骤中，若 \(G\sim\Gamma(N+1,1)\) 且 \(Z=N/G\)，则

\[
\mathbb E(Z-1)=0,\qquad
\mathbb E(Z-1)^2=\frac1{N-1}.
\]

对 \(\Delta(s)=\det(I-s^2B^TB)^{-1/2}\)，固定的算子范数间隙与 Frobenius 范数界保证 \(\Delta''\) 在所用的固定窗口上一致有界。二阶 Taylor 展开和原文的窗口外估计因而给出 \(O(N^{-1})\) 误差。这说明非负双随机情形的永久式误差可以由已有证明细化得到；全局路径估计还需要分别控制缩放、度差和例外顶点造成的误差。

\phantomsection\label{technical-scaling}

\Needspace{12\baselineskip}
\section{4. 对角缩放与非主子矩阵}\label{ux5bf9ux89d2ux7f29ux653eux4e0eux975eux4e3bux5b50ux77e9ux9635}

本节回到图论记号，\(s=S\mathbf1\) 是度差向量，\(a=s/(n-1)\)，并假定 \(n\ge2\)。定理 3.1 适用于经标量归一化后行列和相等的矩阵。竞赛图矩阵通常具有不同的行列和：若 \(C=2A/(n-1)\)，则

\[
C\mathbf1=\mathbf1+a,\qquad C^{\mathsf T}\mathbf1=\mathbf1-a.
\]

因此，修正边际时还须估计永久式及其高斯因子的变化。先构造局部对角缩放，估计其总对数代价；再用 Gram 矩阵比较控制删行、删列的影响；最后，通过成对度差修正，使竞赛图子矩阵满足局部定理的条件，从而证明引理 4.4。

对一般阶数 \(p\ge1\)，沿用 \(P_p=J_p/p=\mathbf1_p\mathbf1_p^{\mathsf T}/p\)，并令 \(\Pi_p=I_p-P_p\)，即去掉全一方向的正交投影。矩阵的总质量为 \(\mathfrak m(X)=\sum_{i,j}X_{ij}\)。双随机矩阵指元素非负、每行和每列之和均为一的方阵。本节常数在给定稠密度与中心核谱隙参数后统一。

\textbf{定义 4.1（边际误差）。}

对任意实 \(p\times p\) 矩阵 \(Q\)（\(p\ge1\)，此处不要求质量已经为 \(p\)），定义

\[
\alpha(Q)=Q\mathbf1_p-\mathbf1_p,\qquad
\beta(Q)=Q^{\mathsf T}\mathbf1_p-\mathbf1_p,
\]

\[
g(Q)=(\alpha(Q),\beta(Q))\in\mathbb R^{2p},\qquad
\varepsilon(Q)=\|g(Q)\|_\infty,
\]

\[
\|g(Q)\|_2^2=\sum_i\alpha(Q)_i^2+\sum_j\beta(Q)_j^2.
\]

\((x,y)\) 同样指把两个向量首尾拼接，故 \(\|(x,y)\|_\infty=\max(\|x\|_\infty,\|y\|_\infty)\)。质量为 \(p\) 时两侧误差的和各为零；质量不为 \(p\) 时，以上定义仍成立，但须先归一化才调用引理 4.1。

以下以 \(n\) 表示全图阶数，以 \(m=n-t\) 表示保留阶数。引理 4.1、4.2 中用 \(p\) 表示一般矩阵阶数，应用时取 \(p=m\)。矩阵 \(C,C^{\prime},\widehat C\) 为 \(n\) 阶，\(X,\widetilde X,B_X,Z\) 为 \(m\) 阶，其总质量分别为

\[
\mathfrak m(C')=n+\nu,\quad \mathfrak m(\widehat C)=n,\quad
\mathfrak m(X)=m+\kappa,\quad
\mathfrak m(\widetilde X)=\mathfrak m(B_X)=m.
\]

这里 \(Z=\Pi_m[(S-I)/(n-1)][R,T]\Pi_m\)，而 \(B_X\) 是缩放后的双随机矩阵。应用定理 3.1 时，扰动矩阵为 \(mB_X-J_m\)，归一化扰动为 \(B_X-P_m\)。

\Needspace{12\baselineskip}
\subsection{4.1 局部缩放及其代价}\label{ux5c40ux90e8ux7f29ux653eux53caux5176ux4ee3ux4ef7}

我们寻找形如 \(B_{ij}=X_{ij}e^{x_i+y_j}\)、行列和均为一的矩阵。将 \((x,y)\) 换为 \((x+c\mathbf1,y-c\mathbf1)\) 不改变 \(B\)；条件 \(\sum x_i=\sum y_j\) 消去这一维自由度。奇异值谱隙控制其余子空间上的线性化平衡方程，而元素界进一步给出不依赖维数的无穷范数控制，从而可以使用压缩映射论证。

\textbf{引理 4.1（局部缩放）。}

设 \(p\ge1\)，\(X\) 为总质量等于 \(p\) 的实 \(p\times p\) 矩阵。记

\[
\alpha=X\mathbf 1_p-\mathbf 1_p,\qquad \beta=X^{\mathsf T}\mathbf 1_p-\mathbf 1_p,\qquad \varepsilon=\max(\|\alpha\|_\infty,\|\beta\|_\infty).
\]

假设 \(|X_{ij}|\le C/p\)，每行和每列的绝对值之和至多为 \(2\)，并且

\[
X_0=X-\frac{\alpha\mathbf 1_p^{\mathsf T}+\mathbf 1_p\beta^{\mathsf T}}p,\qquad E=X_0-P_p
\]

满足 \(|E_{ij}|\le C_0/p\) 和 \(\|E\|_{\rm op}\le q<1\)，其中 \(C,C_0\ge0\)、\(0\le q<1\) 为固定参数。定义

\[
L_0=\frac32+C_0+\frac{C_0^2}{1-q},
\]

\[
\varepsilon_0=\min\left\{\frac1{256L_0^2},\frac{1-q}{2(32L_0+2)}\right\}.
\]

若 \(\varepsilon\le\varepsilon_0\)，则存在满足 \(\sum_i x_i=\sum_j y_j\) 的有限实向量 \(x,y\)，使 \(B=\operatorname{diag}(e^x)X\operatorname{diag}(e^y)\) 的每行和每列之和均为 \(1\)。此外，

\[
\|(x,y)\|_\infty\le4L_0\varepsilon,\qquad |B_{ij}|\le2C/p,\qquad \|B-P_p\|_{\rm op}\le(1+q)/2.
\]

若 \(X\) 逐元素非负，则 \(B\) 为双随机矩阵，且与 \(X\) 具有相同的零支撑。

\textbf{证明。}

两个边际误差向量的元素和均为零。因此，\(X_0\) 的行列和均为 \(1\)，而 \(E\) 双侧中心化。令 \(u=(\mathbf 1_p,-\mathbf 1_p)\)、\(N=uu^{\mathsf T}/(2p)\)。平衡 Hessian \(\mathsf H_0\) 的对角块为 \(I\)，非对角块分别为 \(X_0,X_0^{\mathsf T}\)。其规范化逆可显式写出。置

\[
R_L=(I-EE^{\mathsf T})^{-1},\qquad R_R=(I-E^{\mathsf T}E)^{-1},\qquad O=ER_R.
\]

则 \((\mathsf H_0+N)^{-1}\) 的四个分块为

\[
M_{LL}=R_L-P_p/4,\qquad M_{RR}=R_R-P_p/4,
\]

\[
M_{LR}=-O-P_p/4,\qquad M_{RL}=-O^{\mathsf T}-P_p/4.
\]

这是因为：在共同常数方向上，逆特征值为 \(1/2\)；在规范方向上为 \(1\)；在两个零和子空间上则是标准的分块逆。稠密元素界给出

\[
|(R_L-I)_{ij}|,\ |(R_R-I)_{ij}|\le\frac{C_0^2}{p(1-q^2)},
\]

\[
|O_{ij}|\le\frac{C_0}{p}+\frac{C_0^2q}{p(1-q^2)}.
\]

最后一个估计由分解 \(O=E+EE^{\mathsf T}ER_R\)，再控制两端行、列的二范数得到。向量的一范数为

\[
\|x\|_1=\sum_i|x_i|,
\]

相应的矩阵诱导一范数和无穷范数分别为 \(\|M\|_{1\to1}=\max_j\sum_i|M_{ij}|\)、\(\|M\|_{\infty\to\infty}=\max_i\sum_j|M_{ij}|\)。因此，上述逆矩阵的这两种诱导范数均至多为 \(L_0\)。

令 \(\mathsf H_X\) 的两个对角块分别为 \(\operatorname{diag}(X\mathbf 1_p)\) 和 \(\operatorname{diag}(X^{\mathsf T}\mathbf 1_p)\)，非对角块为 \(X,X^{\mathsf T}\)。差矩阵 \(\mathsf H_X-\mathsf H_0\) 的两种诱导范数均至多为 \(3\varepsilon\)。由 Neumann 级数，\((\mathsf H_X+N)^{-1}\) 的两种诱导范数均至多为 \(2L_0\)。

记 \(z=(x,y)\)、\(g=(\alpha,\beta)\)，平衡方程为

\[
0=g+\mathsf H_Xz+\mathcal N_X(z).
\]

其中，\(\mathcal N_X(z)\) 是 \(X_{ij}(e^{x_i+y_j}-1-x_i-y_j)\) 的行和与列和拼接而成的向量。在球 \(\|z\|_\infty\le R\le1/4\) 上，绝对边际界给出

\[
\|\mathcal N_X(z)\|_\infty\le8R^2,\qquad \|D\mathcal N_X(z)\|_{\infty\to\infty}\le16R.
\]

这些向量均与 \(u\) 正交。在规范子空间 \(u^\perp\) 上考虑

\[
\mathcal T(z)=-(\mathsf H_X+N)^{-1}(g+\mathcal N_X(z)).
\]

取 \(R=4L_0\varepsilon\)。原点的像的范数至多为 \(R/2\)，而 Lipschitz 常数至多为 \(128L_0^2\varepsilon\le1/2\)。因此，映射将规范子空间中的该球映入自身，且为压缩映射，其不动点即为有限势向量。若 \(\varepsilon=0\)，直接取 \(z=0\)。

最后，由 \(R\le1/4\) 得到元素界。行列绝对和的假设给出

\[
\|B-X\|_{1\to1},\ \|B-X\|_{\infty\to\infty}\le32L_0\varepsilon.
\]

又有 \(\|X-X_0\|_{\rm op}\le2\varepsilon\)，由 \(\varepsilon_0\) 定义中的第二个条件即得中心核的谱隙界。正对角因子同时保留非负性与零支撑。引理得证。

仅找到平衡点还不足以估计永久式：恢复对角因子时，需要将永久式乘以 \(\exp(-\sum x_i-\sum y_j)\)。下面的引理既控制这一总代价，也给出比较高斯因子所需的 Frobenius 位移界。

\textbf{引理 4.2（位移与容量）。}

在上述假设下，再设 \(X\ge0\)。以 \(\theta_X=\sum_i x_i+\sum_j y_j\) 表示总缩放势。存在仅依赖固定稠密度与谱隙参数的常数，使

\[
\|(x,y)\|_2\le K\|(\alpha,\beta)\|_2,\qquad \|B-X\|_F\le\frac K{\sqrt p}\|(\alpha,\beta)\|_2,
\]

\[
0\le\theta_X\le K\|(\alpha,\beta)\|_2^2.
\]

\textbf{一般质量。} 对于总质量为任意正数的非负矩阵 \(X\)，假设 \(\widetilde X=pX/\mathfrak m(X)\) 满足引理 4.1 的条件，再对 \(\widetilde X\) 应用以上结论。若 \(B\) 为其缩放，则有精确恒等式

\[
\theta_X=\theta_{\widetilde X}+p\log\frac p{\mathfrak m(X)},\qquad \operatorname{per}X=e^{-\theta_X}\operatorname{per}B.
\]

特别地，若 \(\kappa=\mathfrak m(X)-p\)，则 \(\operatorname{per}X\le e^\kappa\operatorname{per}B\)。当质量不等于 \(p\) 时，\(\theta_X\) 本身不必非负。

\textbf{证明。}

凸势函数为

\[
\Phi_X(x,y)=\sum_{i,j}X_{ij}e^{x_i+y_j}-\sum_i x_i-\sum_j y_j.
\]

其在原点的梯度为 \(g=(\alpha,\beta)\)，平衡点是最小值点。若 \(R=\|(x,y)\|_\infty\)，则沿 \(tz\)、\(0\le t\le1\) 的 Hessian 不小于平衡矩阵 \(B\) 处 Hessian 的 \(e^{-2R}\) 倍。在规范子空间上，后者的最小特征值至少为 \(1-\|B-P_p\|_{\rm op}\)。故整条线段上具有固定正下界 \(\lambda\)。对梯度积分得到 \(\|z\|_2\le\lambda^{-1}\|g\|_2\)。再由稠密元素界和 \(|e^{x_i+y_j}-1|\le K|x_i+y_j|\)，得到 Frobenius 位移估计。

由于 \(\Phi_X(0)=p\)、\(\Phi_X(z)=p-\theta_X\)，凸性给出 \(\theta_X\ge0\)。强凸性给出 \(\theta_X\le\|g\|_2^2/(2\lambda)\)。非单位质量恒等式由将标量 \(p/\mathfrak m(X)\) 吸收入两个对角势中得到。最后，利用 \(p\log(\mathfrak m(X)/p)\le\kappa\) 得到永久式上界。

由质量 \(\sum_{ij}X_{ij}=p\)，

\[
\sum_i\alpha_i=\sum_j\beta_j=0,
\quad X_0\mathbf1=X\mathbf1-\alpha-\mathbf1\frac{\sum_j\beta_j}p=\mathbf1,
\quad X_0^{\mathsf T}\mathbf1=\mathbf1.
\]

因 \(P_p\mathbf1=\mathbf1\)，所以 \(E\mathbf1=E^{\mathsf T}\mathbf1=0\)。同一计算也给出

\[
\Pi_pX\Pi_p=X-P_p-\frac{\alpha\mathbf1^{\mathsf T}+\mathbf1\beta^{\mathsf T}}p=E.
\]

在常数子空间，令 \(c=(\mathbf1,\mathbf1)\)、\(u=(\mathbf1,-\mathbf1)\)，则

\[
\mathsf H_0c=2c,\quad\mathsf H_0u=0,\quad Nc=0,\quad Nu=u.
\]

因此逆在这两条直线上的特征值分别为 \(1/2,1\)。在两个零和子空间上，直接相乘验证

\[
\begin{pmatrix}I&E\\E^{\mathsf T}&I\end{pmatrix}
\begin{pmatrix}R_L&-ER_R\\-E^{\mathsf T}R_L&R_R\end{pmatrix}
=\begin{pmatrix}I&0\\0&I\end{pmatrix},
\]

其中使用

\[
R_L-ER_RE^{\mathsf T}=I,\quad R_R-E^{\mathsf T}R_LE=I,
\quad R_LE=ER_R.
\]

最后一个恒等式可逐项由 Neumann 级数验证。将常数子空间上的投影按原分块表示，得到 \(R_L-P_p/4\) 与 \(-ER_R-P_p/4\) 等四块。

因为每行每列 \(E\) 的二范数至多 \(C_0/\sqrt p\)，

\[
R_L-I=E(I-E^{\mathsf T}E)^{-1}E^{\mathsf T},
\quad |(R_L-I)_{ij}|\le\frac{C_0^2}{p(1-q^2)},
\]

\[
O-E=EE^{\mathsf T}E(I-E^{\mathsf T}E)^{-1},
\quad |(O-E)_{ij}|\le\frac{C_0^2q}{p(1-q^2)}.
\]

每个 \(2p\) 维分块行的绝对和因此至多为

\[
1+\frac14+\frac14+C_0+
\frac{C_0^2(1+q)}{1-q^2}
=\frac32+C_0+\frac{C_0^2}{1-q}=L_0.
\]

矩阵对称，所以诱导一范数相同。又

\[
\|X-X_0\|_{\infty\to\infty},\|X-X_0\|_{1\to1}\le2\varepsilon,
\quad\|\mathsf H_X-\mathsf H_0\|_{\infty\to\infty},
\|\mathsf H_X-\mathsf H_0\|_{1\to1}\le3\varepsilon.
\]

令 \(M=(\mathsf H_0+N)^{-1}\)，\(D=\mathsf H_X-\mathsf H_0\)；则

\[
(\mathsf H_X+N)^{-1}=(I+MD)^{-1}M
=\sum_{h=0}^{\infty}(-MD)^hM,
\]

\[
\|(\mathsf H_X+N)^{-1}\|\le\frac{L_0}{1-3L_0\varepsilon}\le2L_0.
\]

这里任取上述一范数或无穷范数；\(\varepsilon\le1/(256L_0^2)\) 保证 \(3L_0\varepsilon<1/2\)。

对 \(|h|\le2R\le1/2\)，Taylor 积分余项为

\[
e^h-1-h=h^2\int_0^1(1-t)e^{th}\,dt,
\qquad |e^h-1-h|\le\tfrac12e^{1/2}h^2.
\]

每行绝对和至多二、\(|x_i+y_j|\le2R\)，所以

\[
\|\mathcal N_X(z)\|_\infty\le2\cdot\tfrac12e^{1/2}(2R)^2\le8R^2.
\]

导数中对同一行势与所有列势的绝对系数和至多为

\[
2\sum_j|X_{ij}|\,|e^{x_i+y_j}-1|
\le2\cdot2\cdot2Re^{1/2}\le16R.
\]

两侧总和相等，故 \(u^{\mathsf T}g=u^{\mathsf T}\mathcal N_X(z)=0\)；又 \(\mathsf H_Xu=0\)，所以逆保持 \(u^\perp\)。令 \(R=4L_0\varepsilon\)，则

\[
\|\mathcal T(0)\|_\infty\le2L_0\varepsilon=R/2,
\quad\operatorname{Lip}(\mathcal T)\le2L_0\cdot16R
=128L_0^2\varepsilon\le1/2.
\]

因此 \(\|\mathcal T(z)\|_\infty\le R/2+\|z\|_\infty/2\le R\)。从 \(z_0=0\)、\(z_{h+1}=\mathcal T(z_h)\) 出发，

\[
\|z_{h+1}-z_h\|_\infty\le2^{-h}\|z_1-z_0\|_\infty,
\quad\sum_{h=0}^{\infty}\|z_{h+1}-z_h\|_\infty<\infty.
\]

故迭代收敛到有限实向量，极限满足平衡方程。因为 \(R\le1/4\)，\(e^{x_i+y_j}\le e^{1/2}<2\)，从而 \(|B_{ij}|\le2C/p\)。同时

\[
\|B-X\|_{\infty\to\infty},\|B-X\|_{1\to1}
\le2\cdot2Re^{1/2}\le8R=32L_0\varepsilon,
\]

\[
\|B-P_p\|_{\rm op}\le q+(32L_0+2)\varepsilon\le(1+q)/2.
\]

对非负 \(X\)，任意方向 \(h=(v,w)\) 的 Hessian 二次型为

\[
h^{\mathsf T}\nabla^2\Phi_X(z)h
=\sum_{ij}X_{ij}e^{x_i+y_j}(v_i+w_j)^2\ge0.
\]

在平衡点 \(z_*\)，Hessian 为 \(\left(\begin{smallmatrix}I&B\\B^{\mathsf T}&I\end{smallmatrix}\right)\)。在规范子空间上，将常数方向与零和方向分解；前者特征值二，后者下界 \(1-\|B-P_p\|_{\rm op}\)。若 \(0\le t\le1\)，

\[
X_{ij}e^{t(x_i+y_j)}=B_{ij}e^{-(1-t)(x_i+y_j)}\ge e^{-2R}B_{ij}.
\]

故可取 \(\lambda=e^{-2R}(1-\|B-P_p\|_{\rm op})>0\)。由梯度线段积分，

\[
0=\nabla\Phi_X(z_*)=g+\int_0^1\nabla^2\Phi_X(tz_*)z_*\,dt,
\]

\[
\lambda\|z_*\|_2^2\le-\langle g,z_*\rangle\le\|g\|_2\|z_*\|_2,
\quad\|z_*\|_2\le\lambda^{-1}\|g\|_2.
\]

对于 Frobenius 位移，

\[
\begin{aligned}
\|B-X\|_F^2
&=\sum_{ij}X_{ij}^2(e^{x_i+y_j}-1)^2\\
&\le K\frac{C^2}{p^2}\sum_{ij}(x_i+y_j)^2\\
&\le K\frac{C^2}{p^2}\left(2p\sum_ix_i^2+2p\sum_jy_j^2\right)
\le\frac{K'}p\|g\|_2^2.
\end{aligned}
\]

平衡时 \(\sum_{ij}B_{ij}=p\)，所以 \(\Phi_X(z_*)=p-\theta_X\)。强凸性给出

\[
p-\theta_X\ge p+\langle g,z_*\rangle+\frac\lambda2\|z_*\|_2^2,
\]

\[
0\le\theta_X\le-\langle g,z_*\rangle-\frac\lambda2\|z_*\|_2^2
\le\frac{\|g\|_2^2}{2\lambda}.
\]

最后一个不等式由 \(0\le\|\sqrt\lambda z_*+g/\sqrt\lambda\|_2^2\) 得到。若质量为 \(p+\kappa\)，令 \(c=p/(p+\kappa)\)；将 \(\tfrac12\log c\) 同时加入两侧势，得到

\[
\theta_X=\theta_{\widetilde X}+p\log c,\qquad
\operatorname{per}X=e^{-\theta_X}\operatorname{per}B,
\]

\[
-\theta_X\le p\log(1+\kappa/p)\le\kappa.
\]

此处 \(\kappa\) 允许为负，但必须满足 \(p+\kappa>0\)。

为缩短本段公式，暂记 \(P=P_p\)、\(L=R_L\)、\(R=R_R\)、\(O=ER_R\)。中心化给出

\[
EP=PE=E^{\mathsf T}P=PE^{\mathsf T}=0,\quad
LP=PL=RP=PR=P,\quad OP=PO=O^{\mathsf T}P=PO^{\mathsf T}=0.
\]

加入规范投影后的矩阵为

\[
\mathsf H_0+N=\begin{pmatrix}I+P/2&E+P/2\\E^{\mathsf T}+P/2&I+P/2\end{pmatrix},
\]

候选逆为 \(M=\left(\begin{smallmatrix}L-P/4&-O-P/4\\-O^{\mathsf T}-P/4&R-P/4\end{smallmatrix}\right)\)。左上分块逐项相乘：

\[
\begin{aligned}
[(\mathsf H_0+N)M]_{LL}
={}&(I+P/2)(L-P/4)+(E+P/2)(-O^{\mathsf T}-P/4)\\
={}&L-P/4+P/2-P/8-EO^{\mathsf T}-P/8\\
={}&L-EO^{\mathsf T}.
\end{aligned}
\]

因为 \(O^{\mathsf T}=E^{\mathsf T}L\)，

\[
L-EO^{\mathsf T}=L-EE^{\mathsf T}L=(I-EE^{\mathsf T})L=I.
\]

右上分块为

\[
\begin{aligned}
[(\mathsf H_0+N)M]_{LR}
={}&(I+P/2)(-O-P/4)+(E+P/2)(R-P/4)\\
={}&-O-P/4-P/8+ER+P/2-P/8\\
={}&-O+ER=0.
\end{aligned}
\]

左下分块为

\[
\begin{aligned}
[(\mathsf H_0+N)M]_{RL}
={}&(E^{\mathsf T}+P/2)(L-P/4)+(I+P/2)(-O^{\mathsf T}-P/4)\\
={}&E^{\mathsf T}L+P/2-P/8-O^{\mathsf T}-P/4-P/8\\
={}&E^{\mathsf T}L-O^{\mathsf T}=0.
\end{aligned}
\]

右下分块为

\[
\begin{aligned}
[(\mathsf H_0+N)M]_{RR}
={}&(E^{\mathsf T}+P/2)(-O-P/4)+(I+P/2)(R-P/4)\\
={}&-E^{\mathsf T}O-P/8+R-P/4+P/2-P/8\\
={}&R-E^{\mathsf T}O\\
={}&R-E^{\mathsf T}ER=(I-E^{\mathsf T}E)R=I.
\end{aligned}
\]

因此 \(M\) 是所求逆矩阵。

行侧余项和列侧余项分别为

\[
\mathcal N_{L,i}(x,y)=\sum_jX_{ij}(e^{x_i+y_j}-1-x_i-y_j),
\]

\[
\mathcal N_{R,j}(x,y)=\sum_iX_{ij}(e^{x_i+y_j}-1-x_i-y_j).
\]

其行侧偏导数为

\[
\frac{\partial\mathcal N_{L,i}}{\partial x_k}
=\mathbf1\{i=k\}\sum_jX_{ij}(e^{x_i+y_j}-1),\quad
\frac{\partial\mathcal N_{L,i}}{\partial y_j}=X_{ij}(e^{x_i+y_j}-1).
\]

列侧偏导数为

\[
\frac{\partial\mathcal N_{R,j}}{\partial x_i}=X_{ij}(e^{x_i+y_j}-1),\quad
\frac{\partial\mathcal N_{R,j}}{\partial y_k}
=\mathbf1\{j=k\}\sum_iX_{ij}(e^{x_i+y_j}-1).
\]

所以每个行侧 Jacobian 行的绝对和至多为

\[
\left|\sum_jX_{ij}(e^{x_i+y_j}-1)\right|
+\sum_j|X_{ij}||e^{x_i+y_j}-1|
\le2\sum_j|X_{ij}||e^{x_i+y_j}-1|.
\]

当 \(\|(x,y)\|_\infty\le R\le1/4\) 时，右侧至多

\[
2\sum_j|X_{ij}|\cdot2Re^{2R}\le8Re^{1/2}\le16R.
\]

列侧用列绝对和不超过二得到同一个界。用向量值函数沿连接两点的线段积分，

\[
\mathcal N(z)-\mathcal N(w)=\int_0^1D\mathcal N(w+t(z-w))(z-w)\,dt,
\]

所以其 Lipschitz 常数至多 \(16R\)。再乘规范逆的 \(2L_0\)，得到压缩映射常数 \(32L_0R=128L_0^2\varepsilon\)。

令 \(E_B=B-P_p\)，其中 \(P_p=J_p/p\)。对 \(h=(v,w)\in u^\perp\)，有 \(\sum_iv_i=\sum_jw_j\)，所以可以唯一写成

\[
v=c\mathbf1+v_0,\quad w=c\mathbf1+w_0,\quad
\mathbf1^{\mathsf T}v_0=\mathbf1^{\mathsf T}w_0=0.
\]

其范数平方为

\[
\|h\|_2^2=2pc^2+\|v_0\|_2^2+\|w_0\|_2^2.
\]

因为 \(B\mathbf1=B^{\mathsf T}\mathbf1=\mathbf1\)，交叉常数项为零，逐项计算

\[
\begin{aligned}
h^{\mathsf T}\begin{pmatrix}I&B\\B^{\mathsf T}&I\end{pmatrix}h
&=\|v\|^2+\|w\|^2+2v^{\mathsf T}Bw\\
&=4pc^2+\|v_0\|^2+\|w_0\|^2+2v_0^{\mathsf T}E_Bw_0\\
&\ge4pc^2+\|v_0\|^2+\|w_0\|^2-2q_B\|v_0\|\|w_0\|\\
&\ge4pc^2+(1-q_B)(\|v_0\|^2+\|w_0\|^2)\\
&\ge(1-q_B)\|h\|^2,
\end{aligned}
\]

其中 \(q_B=\|E_B\|_{\rm op}<1\)，倒数第二步用 \(2ab\le a^2+b^2\)。这给出位移与容量引理使用的强凸性常数。

\Needspace{12\baselineskip}
\subsection{4.2 高斯删除与中心化}\label{ux9ad8ux65afux5220ux9664ux4e0eux4e2dux5fc3ux5316}

从稠密归一化矩阵中删除 \(t\) 行，删去的行平方范数总量为 \(t/n\) 量级。将对数行列式导数用于对应的 Gram 矩阵损失，可以保留这个量级。随后改用另一侧 Gram 矩阵处理删列。这样，即使保留的是非主子矩阵，也能得到线性的删除误差。

\textbf{定义 4.2（高斯行列式因子）。}

对满足 \(\|Z\|_{\rm op}<1\) 的实矩阵 \(Z\)，定义

\[
\mathcal G(Z)=\det(I-Z^{\mathsf T}Z)^{-1/2}.
\]

改用 \(ZZ^{\mathsf T}\) 得到相同数值；补零行或零列也不改变此值。

\textbf{引理 4.3（Gram 删除与稳定性）。}

设 \(\|Z\|_{\rm op}\le q_*<1\)。删除行集 \(I\) 和列集 \(J\)，剩余集合为 \(R,T\)，则

\[
0\le\log\mathcal G(Z)-\log\mathcal G(Z[R,T])\le\frac{\sum_{i\in I}\|Z_{i,\cdot}\|_2^2+\sum_{j\in J}\|Z_{\cdot,j}\|_2^2}{2(1-q_*^2)}.
\]

\textbf{双侧中心化。} 对阶数为 \(m\ge1\) 且满足 \(\|W\|_{\rm op}\le q_*\) 的方形余矩阵 \(W\)，置 \(u=\mathbf 1_m/\sqrt m\)，则

\[
0\le\log\mathcal G(W)-\log\mathcal G(\Pi_mW\Pi_m)\le\frac{\|W^{\mathsf T}u\|_2^2+\|Wu\|_2^2}{2(1-q_*^2)}.
\]

\textbf{稳定性。} 对于维数相同、算子范数均至多为 \(q_*\) 的两个实矩阵，另有

\[
|\log\mathcal G(U)-\log\mathcal G(V)|\le\frac{\|U\|_F+\|V\|_F}{2(1-q_*^2)}\|U-V\|_F.
\]

\textbf{证明。}

令 \(P_R\) 为投影到保留行坐标的对角投影矩阵。删除行使列侧 Gram 矩阵减少半正定矩阵 \(Z^{\mathsf T}(I-P_R)Z\)，其迹等于被删行的平方范数之和。对 \(F(H)=-\tfrac12\log\det(I-H)\) 和半正定方向 \(D\)，

\[
DF(H)[D]=\tfrac12\operatorname{tr}((I-H)^{-1}D)\le\frac{\operatorname{tr}D}{2(1-q_*^2)}.
\]

积分此界，再改用行侧 Gram 矩阵处理删列。分别投影掉两个全一方向，以同一方法得到中心化界。对于最后一个结论，沿两个 Gram 矩阵之间的线段积分，并使用

\[
\|U^{\mathsf T}U-V^{\mathsf T}V\|_*\le(\|U\|_F+\|V\|_F)\|U-V\|_F,
\]

其中 \(\|\cdot\|_*\) 为迹范数。两个 Gram 矩阵之间的线段仍满足同一算子谱隙条件。

现在设 \(S\) 为 \(n\ge4\) 阶竞赛符号矩阵，\(s=S\mathbf 1_n\)、\(\tau=\|s\|_2^2/(n-1)^2\)。置

\[
Y=\frac{S-I}{n-1},\qquad D_n(S)=\det(I-S^{\mathsf T}S/n^2)^{-1/2}.
\]

反对称性与奇异值成对出现给出

\[
\|Y\|_{\rm op}^2\le q_Y^2=\frac{n(n-1)/2+1}{(n-1)^2}<1.
\]

取整数 \(0\le t<n\)，任删 \(t\) 行与 \(t\) 列，记 \(m=n-t\ge1\)、\(Z=\Pi_mY[R,T]\Pi_m\)。\(Y\) 的每行和每列的平方二范数均为 \(n/(n-1)^2\)。此外，

\[
\|Y\mathbf 1_n\|_2^2=\|Y^{\mathsf T}\mathbf 1_n\|_2^2=\tau+\frac n{(n-1)^2},
\]

这里 \(\mathbf1_I,\mathbf1_J\in\mathbb R^n\) 是原顶点全集 \([n]\) 中集合 \(I,J\) 的指示列向量：坐标属于相应集合时取一，否则取零。

\[
\|Y\mathbf 1_J\|_2^2,\ \|Y^{\mathsf T}\mathbf 1_I\|_2^2\le\frac{t^2n}{(n-1)^2}.
\]

因此，前一引理给出

\[
0\le\log\mathcal G(Y)-\log\mathcal G(Z)\le\frac{t n/(n-1)^2+2[\tau+n(1+t^2)/(n-1)^2]/m}{1-q_Y^2}.
\]

全矩阵的 Gram 恒等式为 \(Y^{\mathsf T}Y=(S^{\mathsf T}S+I)/(n-1)^2\)。它与 \(S^{\mathsf T}S/n^2\) 的差半正定，且其迹为

\[
\vartheta_n=\frac{2n-1}{n(n-1)}+\frac n{(n-1)^2}.
\]

从而 \(0\le\log\mathcal G(Y)-\log D_n(S)\le43/(4n)\)。若 \(m\ge n/2\)，第一个损失至多为 \((24t+18\tau+8)/n\)。两者都是相对于同一全矩阵核的非负损失，故

\[
|\log\mathcal G(Z)-\log D_n(S)|\le\frac{24t+18\tau+11}{n}.
\]

将 \(1-q_Y^2=n(n-3)/(2(n-1)^2)\) 代入，再利用 \(t\le n/2\) 即得上述常数。此估计适用于独立选择的行、列删集 \(I,J\)；反对称性仅用于原全矩阵 \(S\)。

对 \(0\le H\le q_*^2I\)，从谱分解或行列式微分公式得到

\[
\frac{d}{dt}\left[-\frac12\log\det(I-H-tD)\right]
=\frac12\operatorname{tr}((I-H-tD)^{-1}D).
\]

若 \(D\ge0\) 且整条线段仍不超过 \(q_*^2I\)，则右侧在 \([0,\operatorname{tr}D/(2(1-q_*^2))]\) 内。删除行时

\[
H_0=Z^{\mathsf T}P_RZ,\quad
D=Z^{\mathsf T}(I-P_R)Z\ge0,\quad
\operatorname{tr}D=\sum_{i\in I}\sum_jZ_{ij}^2.
\]

删除列时改用 \((P_RZ)(P_RZ)^{\mathsf T}\)；删去的列平方范数不超过原列平方范数，所以两次损失可以相加。中心化使用投影 \(I-uu^{\mathsf T}\)，对应的第一步 Gram 损失为 \(W^{\mathsf T}uu^{\mathsf T}W\)，迹为 \(\|W^{\mathsf T}u\|^2\)；另一步至多为 \(\|Wu\|^2\)。

对任意 \(U,V\)，

\[
U^{\mathsf T}U-V^{\mathsf T}V
=U^{\mathsf T}(U-V)+(U-V)^{\mathsf T}V,
\]

\[
\|U^{\mathsf T}U-V^{\mathsf T}V\|_*
\le\|U\|_F\|U-V\|_F+\|U-V\|_F\|V\|_F.
\]

在 Gram 线段 \(H_t=(1-t)V^{\mathsf T}V+tU^{\mathsf T}U\) 上仍有 \(0\le H_t\le q_*^2I\)。利用 \(|\operatorname{tr}(AB)|\le\|A\|_{\rm op}\|B\|_*\) 再积分，得到引理中的 Lipschitz 常数。

由于 \(S^{\mathsf T}=-S\)，

\[
(S-I)^{\mathsf T}(S-I)=(-S-I)(S-I)=-S^2+I=S^{\mathsf T}S+I.
\]

因此

\[
q_Y^2=\frac{n(n-1)/2+1}{(n-1)^2},\qquad
1-q_Y^2=\frac{n(n-3)}{2(n-1)^2}>0\quad(n\ge4).
\]

每行恰有 \(n\) 个模为 \(1/(n-1)\) 的元素，平方范数为 \(n/(n-1)^2\)。又

\[
Y\mathbf1=(s-\mathbf1)/(n-1),\quad
Y^{\mathsf T}\mathbf1=(-s-\mathbf1)/(n-1),
\]

\[
\|s\mp\mathbf1\|^2=\|s\|^2+n\mp2\sum_i s_i=\|s\|^2+n.
\]

对删列指示向量，逐行有 \(|(Y\mathbf1_J)_i|\le t/(n-1)\)，故平方范数至多 \(nt^2/(n-1)^2\)。保留矩阵 \(W=Y[R,T]\) 的全一方向满足

\[
\begin{aligned}
\|W\mathbf1_m/\sqrt m\|^2
&\le\frac1m\|Y(\mathbf1_n-\mathbf1_J)\|^2\\
&\le\frac2m\left[\tau+\frac n{(n-1)^2}+\frac{nt^2}{(n-1)^2}\right].
\end{aligned}
\]

另一个方向同理。代入删除和中心化界，得到上述损失 \(L\)。把它乘以 \(n\) 并代入谱隙：

\[
nL\le\frac{2nt}{n-3}
+\frac{4(n-1)^2\tau}{m(n-3)}
+\frac{4n(1+t^2)}{m(n-3)}.
\]

当 \(m=n-t\ge n/2\) 时，

\[
\frac n{n-3}\le4,\quad
\frac{4(n-1)^2}{m(n-3)}\le18,\quad
\frac{4n}{m(n-3)}\le8,
\]

\[
\frac{4nt^2}{m(n-3)}\le16t.
\]

最后一个不等式由 \(t/m\le1\) 与 \(n/(n-3)\le4\) 得到。因此 \(nL\le24t+18\tau+8\)。

全阶 Gram 的差为

\[
\left(\frac1{(n-1)^2}-\frac1{n^2}\right)S^{\mathsf T}S
+\frac I{(n-1)^2}\ge0,
\]

其迹为 \(\vartheta_n\)。对数损失至多

\[
\frac{\vartheta_n}{2(1-q_Y^2)}
=\frac{3n^2-3n+1}{n^2(n-3)}\le\frac{43}{4n}.
\]

最后一个不等式等价于 \(31n^2-117n-4\ge0\)；在 \(n=4\) 时为 \(24>0\)，之后递增。若 \(a,b\ge0\)，则 \(|a-b|\le\max(a,b)\)；两项损失均相对 \(\mathcal G(Y)\)，因此

\[
|\log\mathcal G(Z)-\log D_n(S)|
\le\max\left(\frac{24t+18\tau+8}n,\frac{43}{4n}\right)
\le\frac{24t+18\tau+11}n.
\]

\Needspace{12\baselineskip}
\subsection{4.3 竞赛图子矩阵的预处理}\label{ux7adeux8d5bux56feux5b50ux77e9ux9635ux7684ux9884ux5904ux7406}

矩阵 \(C=2A/(n-1)\) 的行、列误差符号相反。将第 \(i\) 行乘以 \((1+a_i)^{-1}\) 可修正原行和，将第 \(i\) 列乘以 \((1-a_i)^{-1}\) 可修正原列和。两种修正同时进行后，边际仍有误差，以下估计这一误差。其主要作用是使总质量的一阶变化消失，并在恢复乘积中产生度差修正因子 \(\Gamma\)。

固定 \(0\le a_0<1\) 和有限正参数 \(A_0,B_0\)。本小节假设

\[
a_i=\frac{s_i}{n-1},\qquad \max_i|a_i|\le a_0,\qquad \tau=\sum_i a_i^2\le A_0\log n,\qquad t\le B_0\log n.
\]

定义

\[
C=\frac{J-I+S}{n-1}=\frac{2A}{n-1},\qquad \ell_i=\frac1{1+a_i},\qquad r_i=\frac1{1-a_i},\qquad g_i=\frac1{1-a_i^2}.
\]

置 \(C'=\operatorname{diag}(\ell)C\operatorname{diag}(r)\)、\(F=C'-C\)、\(\nu=\mathfrak m(C')-n\)。为计算总质量，记 \(v_i=a_i/(1-a_i^2)\)、\(w_i=a_i^2/(1-a_i^2)\)，以及 \(V=\sum_i v_i\)、\(W=\sum_i w_i\)。展开 \(\ell=\mathbf 1-v+w\) 和 \(r=\mathbf 1+v+w\)，得到精确恒等式

\[
\nu=\frac{W^2-V^2+W+2w^{\mathsf T}Sv}{n-1}.
\]

确切地说，\(\mathbf 1^{\mathsf T}Sv=-(n-1)W\) 抵消了一阶质量变化。固定度差界给出 \(W\le K\tau\)、\(|V|\le K\tau\)、\(\|v\|_2\le K\sqrt\tau\)。逐行估计 \(Sv\) 得到

\[
|\nu|\le K\mathcal M_\tau,\qquad \mathcal M_\tau=\frac{\tau+\tau^2}{n}+\frac{\tau^{3/2}}{\sqrt n}=o(1).
\]

边际恒等式为

\[
C'\mathbf 1_n-\mathbf 1_n=\operatorname{diag}(\ell)C(r-\mathbf 1_n),
\]

\[
(C')^{\mathsf T}\mathbf 1_n-\mathbf 1_n=\operatorname{diag}(r)C^{\mathsf T}(\ell-\mathbf 1_n).
\]

由于 \(C\) 的每行和每列二范数为 \(O(n^{-1/2})\)，且 \(\|C\|_{\rm op}=O(1)\)，无穷范数边际误差为 \(O(\sqrt{\tau/n})\)，二范数边际误差为 \(O(\sqrt\tau)\)。直接逐元素估计还得到

\[
\|F\|_F\le K\sqrt{\tau/n}.
\]

归一化 \(\widehat C=nC'/\mathfrak m(C')\)。以 \(e_i,f_j\) 表示其行、列误差，故 \(\sum e_i=\sum f_j=0\)。对大小均为 \(t\) 的任意删集 \(I,J\)，记 \(R=[n]\setminus I\)、\(T=[n]\setminus J\)、\(m=n-t\)，并令

\[
X=\frac n m\widehat C[R,T],\qquad \kappa=\mathfrak m(X)-m,\qquad \widetilde X=\frac m{\mathfrak m(X)}X.
\]

精确删行列质量恒等式为

\begin{equation}
\kappa=\frac n m\left[-\sum_{i\in I}e_i-\sum_{j\in J}f_j+\sum_{i\in I,j\in J}\widehat C_{ij}-\frac{t^2}{n}\right].
\tag{4.1}
\label{eq-deletion-mass}
\end{equation}

每个保留行在最后一次归一化前的边际误差为

\[
(X\mathbf 1_m)_i-1=\frac n m\left[e_i+\frac t n-\sum_{j\in J}\widehat C_{ij}\right],
\]

列侧有相同形式的对应恒等式。因此，

\[
|\kappa|\le K\left[t\sqrt{\tau/n}+t^2/n\right],
\]

\[
\varepsilon(\widetilde X)\le K(\sqrt{\tau/n}+t/n),\qquad \|g(\widetilde X)\|_2\le K(\sqrt\tau+t/\sqrt n).
\]

这里 \(g(\widetilde X)\) 为行、列边际误差的拼接向量。\(\widetilde X\) 的元素均在 \([0,K/m]\) 内。

为应用引理 4.1，下面验证 \(\widetilde X\) 的中心核条件。记

\[
\eta=\frac{n^2}{\mathfrak m(C')\mathfrak m(X)}=1+O((t+1)/n),
\]

则

\[
\Pi_m\widetilde X\Pi_m=\eta\left[Z+\Pi_mF[R,T]\Pi_m\right].
\]

由 \(\mathfrak m(C^{\prime})=n+o(1)\) 和 \(\mathfrak m(X)=m+o(1)\)，有 \(\eta=(n/m)(1+o(1))\)。压缩不增大 \(Y\) 的算子范数，第二项则由上述 Frobenius 界控制，故中心核的算子范数至多为 \(1/\sqrt2+o(1)\)。其元素为 \(O(1/m)\)，而原矩阵 \(\widetilde X\) 的边际误差无穷范数趋于零。又因 \(\widetilde X\) 的质量为 \(m\)，引理 4.1 中的 \(E\) 等于 \(\Pi_m\widetilde X\Pi_m\)。因此可固定 \(q=3/4\) 及元素界 \(C,C_0\)。相应的 \(L_0,\varepsilon_0\) 仅依赖 \(a_0,A_0,B_0\)，所以引理 4.1 对充分大的 \(n\) 一致适用，给出双随机矩阵 \(B_X\) 及谱范数界 \(7/8\)。

二范数位移引理给出

\[
\|B_X-\widetilde X\|_F\le K(\sqrt{\tau/n}+t/n).
\]

与较强的中心核界结合，\(B_X\) 的中心算子范数仍至多为 \(1/\sqrt2+o(1)\)，例如最终至多为 \(4/5\)。元素继续满足 \(K/m\) 的界。

对 \(Z\) 应用高斯删除估计，并以对数 Lipschitz 界估计预缩放、归一化及对角缩放产生的扰动。由 \(B_X-P_m=\Pi_mB_X\Pi_m\)，得到

\[
|\log\mathcal G(B_X-P_m)-\log D_n(S)|\le K\left[\sqrt{\tau/n}+(t+1)/n\right].
\]

中心化产生的 \(\tau/n\) 项已吸收入 \(\sqrt{\tau/n}\)。其中由删除产生的误差与 \(t/n\) 成线性关系，正适合对后续子集权重求和。

综上，各项误差如下。所有常数与最终适用阈值只依赖固定的 \(a_0,A_0,B_0\)。全阶质量误差在恢复对数中产生 \(O(|\nu|)\le K\mathcal M_\tau\)；删除引起的质量变化至多为 \(K[t\sqrt{\tau/n}+t^2/n]\)；高斯对数误差为 \(K[\sqrt{\tau/n}+(t+1)/n]\)。在阶数 \(m\ge n/2\) 使用定理 3.1，增加 \(O(1/m)=O(1/n)\)。恢复标量因子得到 \(e^{-1}\exp(O((t+1)/n))\)。对于质量为 \(m+\kappa\) 的 \(X\)，容量仅以 \(e^{-\theta_X}\le e^\kappa\) 的形式使用，不能断言这种非单位质量下的 \(\theta_X\) 非负。第 4.4 节精确恢复这些因子，得到引理 4.4 的误差指数。

逐元素验证

\[
1-v_i+w_i=1-\frac{a_i}{1-a_i^2}+\frac{a_i^2}{1-a_i^2}
=\frac1{1+a_i}=\ell_i,
\]

\[
1+v_i+w_i=\frac1{1-a_i}=r_i,\qquad
a_iv_i=w_i,\quad \ell_ir_i=g_i=1+w_i.
\]

令 \(L=\mathbf1-v+w\)、\(R=\mathbf1+v+w\)（这里 \(L,R\) 仅在本段表示向量）。则

\[
L^{\mathsf T}JR=(n-V+W)(n+V+W)=(n+W)^2-V^2,
\]

\[
L^{\mathsf T}IR=\sum_i g_i=n+W.
\]

反对称部分逐项展开为

\[
\begin{aligned}
L^{\mathsf T}SR
={}&\mathbf1^{\mathsf T}Sv+\mathbf1^{\mathsf T}Sw
-v^{\mathsf T}S\mathbf1-v^{\mathsf T}Sv-v^{\mathsf T}Sw\\
&+w^{\mathsf T}S\mathbf1+w^{\mathsf T}Sv+w^{\mathsf T}Sw\\
={}&-2(n-1)W+2w^{\mathsf T}Sv.
\end{aligned}
\]

这里 \(\mathbf1^{\mathsf T}Sv=-(n-1)a^{\mathsf T}v=-(n-1)W\)；\(\mathbf1^{\mathsf T}Sw+w^{\mathsf T}S\mathbf1=0\)；\(v^{\mathsf T}Sv=w^{\mathsf T}Sw=0\)；\(-v^{\mathsf T}Sw=w^{\mathsf T}Sv\)。所以

\[
\begin{aligned}
\mathfrak m(C')&=\frac{(n+W)^2-V^2-(n+W)-2(n-1)W+2w^{\mathsf T}Sv}{n-1}\\
&=n+\frac{W^2-V^2+W+2w^{\mathsf T}Sv}{n-1}.
\end{aligned}
\]

由 \(\sum_i a_i=0\)，

\[
V=\sum_i\left(a_i+\frac{a_i^3}{1-a_i^2}\right)
=\sum_i\frac{a_i^3}{1-a_i^2},\quad
|V|\le\frac{a_0}{1-a_0^2}\tau,
\quad W\le\frac\tau{1-a_0^2},
\]

\[
\|v\|_2\le\frac{\sqrt\tau}{1-a_0^2},\quad
\|Sv\|_\infty\le\sqrt n\|v\|_2,\quad
|w^{\mathsf T}Sv|\le\|w\|_1\|Sv\|_\infty\le K\sqrt n\tau^{3/2}.
\]

代回得到 \(|\nu|\le K[(\tau+\tau^2)/n+\tau^{3/2}/\sqrt n]\)。

因 \(\operatorname{diag}(\ell)C\mathbf1=\mathbf1\)，

\[
C'\mathbf1-\mathbf1=\operatorname{diag}(\ell)C(r-\mathbf1).
\]

且 \(|r_i-1|\le K|a_i|\)、\(|\ell_i-1|\le K|a_i|\)、\(\|C_{i,\cdot}\|_2\le K/\sqrt n\)。所以

\[
\|C'\mathbf1-\mathbf1\|_\infty\le K\sqrt{\tau/n},\quad
\|C'\mathbf1-\mathbf1\|_2\le K\|C\|_{\rm op}\sqrt\tau\le K'\sqrt\tau.
\]

列侧同理。逐元素

\[
F_{ij}=C_{ij}(\ell_ir_j-1),\quad
|\ell_ir_j-1|\le K(|a_i|+|a_j|),
\]

\[
\|F\|_F^2\le\frac K{n^2}\sum_{ij}(a_i^2+a_j^2)
=\frac{2K\tau}n.
\]

对质量为 \(n\) 的 \(\widehat C\)，保留块的质量由容斥直接给出

\[
\begin{aligned}
\sum_{i\in R,j\in T}\widehat C_{ij}
&=n-\sum_{i\in I}(1+e_i)-\sum_{j\in J}(1+f_j)
+\sum_{i\in I,j\in J}\widehat C_{ij}\\
&=n-2t-\sum_Ie_i-\sum_Jf_j+\sum_{I,J}\widehat C_{ij}.
\end{aligned}
\]

而 \(m^2/n=n-2t+t^2/n\)，所以乘 \(n/m\) 再减 \(m\) 即得式 (4.1)。由 \(\widehat C_{ij}\le K/n\)，

\[
|\kappa|\le\frac nm\left(2t\max(\|e\|_\infty,\|f\|_\infty)
+(K+1)t^2/n\right)
\le K'[t\sqrt{\tau/n}+t^2/n].
\]

保留行的误差同样逐项为

\[
\frac nm\left(1+e_i-\sum_{j\in J}\widehat C_{ij}\right)-1
=\frac nm\left(e_i+t/n-\sum_{j\in J}\widehat C_{ij}\right).
\]

归一化乘数 \(m/(m+\kappa)\) 与一相差 \(O(|\kappa|/m)\)；故边际的无穷范数和二范数分别为 \(K(\sqrt{\tau/n}+t/n)\)、\(K(\sqrt\tau+t/\sqrt n)\)。

从定义依次代入，

\[
\widetilde X=\frac m{m+\kappa}\frac nm\frac n{n+\nu}C'[R,T]
=\frac{n^2}{(n+\nu)(m+\kappa)}C'[R,T]=\eta C'[R,T].
\]

又 \(C=J/(n-1)+Y\)，左右乘 \(\Pi_m\) 消掉 \(J\)，得到中心核恒等式。由于 \(m=n-t\)，

\[
\eta=\frac nm(1+\nu/n)^{-1}(1+\kappa/m)^{-1}
=1+O((t+1)/n).
\]

因此

\[
\|\Pi_m\widetilde X\Pi_m-Z\|_F
\le|\eta-1|\|Z\|_F+|\eta|\|F[R,T]\|_F
\le K[(t+1)/n+\sqrt{\tau/n}].
\]

加上对角缩放的位移，

\[
\|(B_X-P_m)-Z\|_F
\le K[(t+1)/n+\sqrt{\tau/n}].
\]

上述矩阵 Frobenius 范数均为 \(O(1)\)，算子范数最终至多固定的 \(4/5<1\)。Gram Lipschitz 界与删除界相加得

\[
\begin{aligned}
|\log\mathcal G(B_X-P_m)-\log D_n(S)|
&\le K\left[\frac{t+\tau+1}n+\sqrt{\tau/n}\right]\\
&\le K'\left[\frac{t+1}n+\sqrt{\tau/n}\right],
\end{aligned}
\]

其中充分大 \(n\) 时 \(\tau/n\le1\)，故 \(\tau/n\le\sqrt{\tau/n}\)。

先记

\[
e'=C'\mathbf1-\mathbf1,\quad
f'=(C')^{\mathsf T}\mathbf1-\mathbf1,\quad
\widehat C=\frac n{n+\nu}C'.
\]

逐坐标有

\[
\begin{aligned}
e_i&=(\widehat C\mathbf1)_i-1
=\frac n{n+\nu}(1+e'_i)-1
=\frac{ne'_i-\nu}{n+\nu},\\
f_j&=(\widehat C^{\mathsf T}\mathbf1)_j-1
=\frac n{n+\nu}(1+f'_j)-1
=\frac{nf'_j-\nu}{n+\nu}.
\end{aligned}
\]

因为 \(\sum_ie'_i=\sum_jf'_j=\nu\)，这些公式也验证

\[
\sum_ie_i=\frac{n\nu-n\nu}{n+\nu}=0,\quad
\sum_jf_j=0.
\]

当 \(|\nu|\le n/2\) 时，

\[
\|e\|_\infty\le2\|e'\|_\infty+2|\nu|/n,\quad
\|e\|_2\le2\|e'\|_2+2|\nu|/\sqrt n,
\]

\[
\|f\|_\infty\le2\|f'\|_\infty+2|\nu|/n,\quad
\|f\|_2\le2\|f'\|_2+2|\nu|/\sqrt n.
\]

对 \(0<\tau\le A_0\log n\)，将 \(|\nu|/n\) 除以 \(\sqrt{\tau/n}\) 后，已有质量界给出至多常数乘

\[
\frac{\sqrt\tau+\tau^{3/2}}{n^{3/2}}+\frac\tau n=o(1).
\]

将 \(|\nu|/\sqrt n\) 除以 \(\sqrt\tau\) 得到同一个表达式。\(\tau=0\) 时 \(a=0\)，所以 \(\nu=0\)，无需作除法。这就分别推出两侧的无穷范数界 \(K\sqrt{\tau/n}\) 和二范数界 \(K\sqrt\tau\)。

设

\[
h=X\mathbf1_m-\mathbf1_m,\quad
k=X^{\mathsf T}\mathbf1_m-\mathbf1_m,\quad
\widetilde X=\frac m{m+\kappa}X.
\]

那么归一化后的两侧误差精确为

\[
\widetilde\alpha_i=\frac{mh_i-\kappa}{m+\kappa},\qquad
\widetilde\beta_j=\frac{mk_j-\kappa}{m+\kappa}.
\]

由 \(\sum_ih_i=\sum_jk_j=\kappa\)，再次验证 \(\sum_i\widetilde\alpha_i=\sum_j\widetilde\beta_j=0\)。当 \(|\kappa|\le m/2\)，

\[
\|\widetilde\alpha\|_\infty\le2\|h\|_\infty+2|\kappa|/m,
\quad\|\widetilde\alpha\|_2\le2\|h\|_2+2|\kappa|/\sqrt m,
\]

\[
\|\widetilde\beta\|_\infty\le2\|k\|_\infty+2|\kappa|/m,
\quad\|\widetilde\beta\|_2\le2\|k\|_2+2|\kappa|/\sqrt m.
\]

这里 \(k\) 为列误差向量。保留块的行、列误差分别为

\[
h_i=\frac nm\left(e_i+t/n-\sum_{j\in J}\widehat C_{ij}\right),\qquad
k_j=\frac nm\left(f_j+t/n-\sum_{i\in I}\widehat C_{ij}\right).
\]

由于 \(\widehat C_{ij}\le K/n\)，

\[
\|h\|_\infty,\|k\|_\infty\le K'[\sqrt{\tau/n}+t/n],
\]

\[
\|h\|_2,\|k\|_2\le K'[\sqrt\tau+t/\sqrt n].
\]

质量项 \(|\kappa|/\sqrt m\le K[t\sqrt\tau/n+t^2/n^{3/2}]\) 被上述二范数界吸收，因为 \(t/n\to0\)；无穷范数质量项也被相应界吸收。局部缩放引理的行、列边际误差条件因而均成立。

\textbf{局部定理的稠密度、谱隙与阈值次序。}

归一化质量为 \(m\)，所以中心矩阵精确为

\[
E=\widetilde X-P_m-\frac{\widetilde\alpha\mathbf1^{\mathsf T}
+\mathbf1\widetilde\beta^{\mathsf T}}m
=\Pi_m\widetilde X\Pi_m.
\]

若 \(0\le\widetilde X_{ij}\le C_1/m\)、两侧边际误差的无穷范数至多 \(\epsilon_n\)，则

\[
|E_{ij}|\le\frac{C_1+1+2\epsilon_n}{m}\le\frac{C_1+3}{m}
\quad(\epsilon_n\le1).
\]

非负性使每个行列绝对和就是行列和，至多 \(1+\epsilon_n\le2\)。而

\[
\|E\|_{\rm op}\le|\eta|\|Z\|_{\rm op}+|\eta|\|F[R,T]\|_F
\le\frac1{\sqrt2}+o(1).
\]

先固定 \(q=3/4\) 和元素界 \(C_0=C_1+3\)；它们确定 \(L_0,\varepsilon_0\)。再选统一阶数阈值，使

\[
\|E\|_{\rm op}\le3/4,\quad
\epsilon_n\le\varepsilon_0,\quad m\ge n/2,
\quad |\nu|\le n/2,\quad |\kappa|\le m/2.
\]

这些条件全由只含固定参数和 \(n\) 的上界保证，所以阈值在选择具体图和删集之前确定。局部定理首先给 \(\|B_X-P_m\|_{\rm op}\le7/8\)；更精细地，

\[
\|B_X-P_m\|_{\rm op}
\le\|\Pi_m\widetilde X\Pi_m\|_{\rm op}
+\|B_X-\widetilde X\|_F
\le1/\sqrt2+o(1)\le4/5.
\]

这一估计同样对竞赛图及删集一致。

取充分大的奇数 \(n\) 及相应的正则圆环竞赛图，记其邻接矩阵为 \(A_n\)、符号矩阵为 \(S_n^{\mathrm{cir}}\)。此时所有度差为零，前面的预缩放因子均为一，而且

\[
C=\frac{2A_n}{n-1}=\frac{J_n-I_n+S_n^{\mathrm{cir}}}{n-1},
\qquad C\mathbf1_n=C^{\mathsf T}\mathbf1_n=\mathbf1_n.
\]

因此 \(C'=X=\widetilde X=C\)（这里不删行列，\(t=0\)）。势向量取 \(x=y=0\) 就已满足平衡方程；局部缩放中的唯一小解给出 \(B_X=C\)。在常数向量方向上 \(C-P_n\) 为零；在其正交补上，

\[
(C-P_n)|_{\mathbf1_n^\perp}
=\left.\frac{S_n^{\mathrm{cir}}-I_n}{n-1}\right|_{\mathbf1_n^\perp},\qquad
(S_n^{\mathrm{cir}}-I_n)^{\mathsf T}(S_n^{\mathrm{cir}}-I_n)=I_n-(S_n^{\mathrm{cir}})^2.
\]

A.2 中的交替切换与置换保留符号矩阵的全部频率；引理 6.1 给出最大正频率 \(\cot(\pi/(2n))\)。因此

\[
\begin{aligned}
\|B_X-P_n\|_{\rm op}
&=\frac{\sqrt{1+\cot^2(\pi/(2n))}}{n-1}\\
&=\frac1{(n-1)\sin(\pi/(2n))}\\
&=\frac{2n}{\pi(n-1)}
\frac{\pi/(2n)}{\sin(\pi/(2n))}
\longrightarrow\frac2\pi<\frac1{\sqrt2}.
\end{aligned}
\]

这里使用了 \(\sin u/u\to1\)。可见中心范数一般不趋于零；固定上界 \(4/5\) 已足以应用定理 3.1。

\Needspace{12\baselineskip}
\subsection{4.4 非主永久式界的证明}\label{ux975eux4e3bux6c38ux4e45ux5f0fux754cux7684ux8bc1ux660e}

下面将局部缩放和删除估计应用于竞赛图子矩阵。这里允许行删集与列删集不同，因为圈覆盖进入、离开异常集时可能使用不同的核心顶点。

\textbf{引理 4.4（非主永久式界）。}

固定 \(0\le b<1\) 及正常数 \(A_0,B_0\)。对 \(N\ge2\) 阶竞赛图核心，置 \(a=S\mathbf1/(N-1)\)，\(\tau=\sum_i a_i^2\)。令 \(I,J\subseteq[N]\) 分别为行、列删集；此处 \(I,J\) 为集合，和单位矩阵、全一矩阵的记号由下文的集合运算区分。假设

\[
\max_i|a_i|\le b,\qquad
\tau\le A_0\log N,\qquad
|I|=|J|=t\le B_0\log N.
\]

定义 \(\ell_i=(1+a_i)^{-1}\)、\(r_i=(1-a_i)^{-1}\)，以及

\[
\Gamma=\prod_i(1-a_i^2),\qquad
\mathcal M_\tau=\frac{\tau+\tau^2}{N}+\frac{\tau^{3/2}}{\sqrt N}.
\]

令 \(R=[N]\setminus I\)、\(T=[N]\setminus J\) 分别为保留行、列集；本式的 \(T\) 是列指标集。则有

\[
\operatorname{per}A[R,T]\le
e^{-1}D_N(S)\Gamma
\left(\prod_{i\in I}\ell_i\right)
\left(\prod_{j\in J}r_j\right)
\frac{(N-t)!}{2^{N-t}}\exp(\varepsilon_{\tau,t}),
\]

\[
\varepsilon_{\tau,t}\le K\left[
\mathcal M_\tau+\sqrt{\tau/N}+\frac{t+1}{N}
+t\sqrt{\tau/N}+\frac{t^2}{N}\right].
\]

对所有充分大的 \(N\)，该估计对每个满足条件的竞赛图及每对 \(I,J\) 同时成立。常数 \(K\) 和起始阶数仅依赖 \(b,A_0,B_0\)。

因子 \(\Gamma\) 记录度数不平衡的代价。对主删集 \(I=J=U\)，恢复的顶点权重为 \(\ell_i r_i=(1-a_i^2)^{-1}\)；删集不同时，行、列因子分别保留。误差中含有 \(\sqrt{\tau/N}\) 量级的项，因此在最后一种情形中，必须保留 \(\Gamma\le e^{-\tau}\)。

证明中以 \(n\) 表示全核心阶数 \(N\)，并取 \(a_0=b\)。前面的构造已给出缩放矩阵 \(B_X\)。现恢复全部对角因子和标量因子，再将所得高斯行列式与 \(D_n(S)\) 比较。

\textbf{证明。}

由前面的构造，\(B_X\) 满足定理 3.1 的元素界和谱隙条件。因此

\[
\operatorname{per}B_X=\frac{m!}{m^m}\mathcal G(B_X-P_m)(1+O(1/m)).
\]

行列缩放因子精确恢复为

\[
\prod_{i\in R}(1+a_i)\prod_{j\in T}(1-a_j)=\Gamma\prod_{i\in I}\ell_i\prod_{j\in J}r_j.
\]

因此，删除行恢复 \(\ell\) 因子，删除列恢复 \(r\) 因子。标量因子满足

\begin{equation}
\operatorname{per}A[R,T]=\Gamma\prod_{i\in I}\ell_i\prod_{j\in J}r_j\left[\frac{m(n-1)\mathfrak m(C')}{2n^2}\right]^m e^{-\theta_X}\operatorname{per}B_X.
\tag{4.2}
\label{eq-exact-restoration}
\end{equation}

将 \(m!/m^m\) 代入\hyperref[eq-exact-restoration]{式 (4.2)}后，剩余标量为

\[
\frac{m!}{2^m}(1-1/n)^m\left[\frac{\mathfrak m(C')}{n}\right]^m.
\]

其对数与 \(e^{-1}m!/2^m\) 的对数之差为 \(O((t+1)/n+|\nu|)\)。容量界给出 \(e^{-\theta_X}\le e^\kappa\)。代入质量、高斯和零阶永久式误差，即可得到所需指数界。以下展开恢复因子的计算，所有常数均对竞赛图和删集一致。

行、列删集的因子分别为

\[
\begin{aligned}
\operatorname{per}C[R,T]
&=\left(\prod_{i\in R}\ell_i^{-1}\right)
\left(\prod_{j\in T}r_j^{-1}\right)\operatorname{per}C'[R,T]\\
&=\left(\prod_{i=1}^n(1+a_i)\prod_{j=1}^n(1-a_j)\right)
\left(\prod_{i\in I}(1+a_i)^{-1}\right)
\left(\prod_{j\in J}(1-a_j)^{-1}\right)\operatorname{per}C'[R,T]\\
&=\Gamma\prod_{i\in I}\ell_i\prod_{j\in J}r_j\operatorname{per}C'[R,T].
\end{aligned}
\]

从 \(X=n^2C'[R,T]/[m(n+\nu)]\)，得到

\[
\operatorname{per}C'[R,T]=[m(n+\nu)/n^2]^m\operatorname{per}X.
\]

再用 \(A=(n-1)C/2\) 与 \(\operatorname{per}X=e^{-\theta_X}\operatorname{per}B_X\)，即得式 (4.2)。代入永久式近似后的标量为

\[
\left[\frac{m(n-1)(n+\nu)}{2n^2}\right]^m\frac{m!}{m^m}
=\frac{m!}{2^m}(1-1/n)^m(1+\nu/n)^m.
\]

令 \(u=1/n\)。因

\[
|\log(1-u)+u|\le\frac{u^2}{2(1-u)},
\]

\[
\left|m\log(1-1/n)+1\right|
\le\left|1-\frac mn\right|+\frac{m}{2n^2(1-1/n)}
\le\frac tn+\frac1{2(n-1)}.
\]

当 \(|\nu|/n\le1/2\)，\(|m\log(1+\nu/n)|\le2|\nu|\)。容量部分满足 \(-\theta_X\le\kappa\)。高斯因子满足 \(\mathcal G(B_X-P_m)\le D_n(S)e^{K[\sqrt{\tau/n}+(t+1)/n]}\)。最后 \(1+O(1/m)\le e^{K/n}\)（增大统一阈值），所以误差指数逐项为

\[
\begin{aligned}
\varepsilon_{\tau,t}\le K\bigg[&\underbrace{\mathcal M_\tau}_{\text{全阶质量}}
+\underbrace{t\sqrt{\tau/n}+t^2/n}_{\text{删除质量}}\\
&+\underbrace{\sqrt{\tau/n}+(t+1)/n}_{\text{高斯与标量}}
+\underbrace{1/n}_{\text{永久式误差}}\bigg].
\end{aligned}
\]

合并相同量级即为引理 4.4。 \(\square\)

\Needspace{12\baselineskip}
\subsection{4.5 小度差时的双侧估计}\label{ux5c0fux5ea6ux5deeux65f6ux7684ux53ccux4fa7ux4f30ux8ba1}

当每个度差均为 \(o(\sqrt n)\) 时，原始边际已足够接近一，可以直接使用局部缩放。容量界随后从两侧控制恢复因子，得到主文定理 8.1 所需的近似。

\textbf{引理 4.5（小度差永久式近似）。}

设 \(\Delta_s=\|S\mathbf 1_n\|_\infty=o(\sqrt n)\)。固定 \(0<B_0<\infty\)，任删 \(t\le B_0\log n\) 行与 \(t\) 列，记 \(m=n-t\)。对这些选择一致地有

\begin{equation}
\operatorname{per}A[R,T]=e^{-1}D_n(S)\frac{m!}{2^m}\left[1+O_{B_0}\left(\frac{(\Delta_s+t+1)^2}{n}\right)\right].
\tag{4.3}
\label{eq-small-score-permanent}
\end{equation}

更具体地，固定 \(B_0>0\) 后，存在常数 \(C_{\mathrm{err}},\delta>0\) 和整数 \(N_0\ge4\)，使下述结论对所有 \(n\ge N_0\) 一致成立。设 \(S\) 为竞赛符号矩阵，\(I,J\subseteq[n]\) 且 \(|I|=|J|=t\le B_0\log n\)；置 \(\Delta_s=\|S\mathbf1_n\|_\infty\)、\(\eta=(\Delta_s+t+1)^2/n\)。若 \(\eta\le\delta\)，则\hyperref[eq-small-score-permanent]{式 (4.3)}的相对误差绝对值不超过 \(C_{\mathrm{err}}\eta\)。这里 \(R=[n]\setminus I\)、\(T=[n]\setminus J\)、\(m=n-t\)，两个删集可独立选取。在引理所述渐近范围内，\(\eta\to0\)。

\textbf{证明。}

增大 \(N_0\)，使整个对数删除窗口内均有 \(2t\le n\)，于是 \(m\ge n/2>0\)。此处不需要成对预缩放。令 \(C=2A/(n-1)\)、\(X=(n/m)C[R,T]\)。全矩阵 \(C\) 的行列误差分别为 \(s_i/(n-1)\)、\(-s_i/(n-1)\)。因此，精确质量恒等式\hyperref[eq-deletion-mass]{式 (4.1)}及第 4.3 节的边际公式给出

\[
|\kappa|=O(t(\Delta_s+t)/n),\qquad \varepsilon(X)=O((\Delta_s+t)/n),\qquad \|g(X)\|_2=O((\Delta_s+t)/\sqrt n).
\]

质量误差满足 \(|\kappa|\le K\eta\)。必要时缩小 \(\delta\)，便一致保证 \(\mathfrak m(X)=m+\kappa>0\)。归一化 \(\widetilde X=mX/\mathfrak m(X)\)。其稠密度与中心核谱隙固定，局部缩放引理适用。二范数位移和容量估计给出

\[
\|B_X-\widetilde X\|_F=O((\Delta_s+t)/n),\qquad 0\le\theta_{\widetilde X}=O((\Delta_s+t)^2/n).
\]

非单位质量修正随后给出 \(|\theta_X|=O((\Delta_s+t)^2/n)\)。又有 \(\tau\le n\Delta_s^2/(n-1)^2\)。有限删除估计、归一化因子和对角缩放的扰动共同给出

\[
|\log\mathcal G(B_X-P_m)-\log D_n(S)|=O((\Delta_s+t+1)^2/n).
\]

此时不含度差乘积的精确恢复为

\[
\operatorname{per}A[R,T]=\left[\frac{m(n-1)}{2n}\right]^m e^{-\theta_X}\operatorname{per}B_X.
\]

应用统一永久式定理，并使用 \((1-1/n)^m=e^{-1}\exp(O((t+1)/n))\)。所有对数误差均由 \(\eta=(\Delta_s+t+1)^2/n\) 的一个固定倍数控制。将 \(\delta\) 取得充分小，对这些界取指数，便在前述有限范围内一致得到相对误差绝对值不超过 \(C_{\mathrm{err}}\eta\)。在 \(\Delta_s=o(\sqrt n)\)、\(t\le B_0\log n\) 的渐近范围中，\(\eta\to0\)。由此得到主文定理 8.1 所用的短删集子矩阵的一致近似。

因 \(|s_i|\le \Delta_s\)，全矩阵边际绝对误差至多 \(\Delta_s/(n-1)\)，删集质量公式给出

\[
|\kappa|\le\frac nm\left[\frac{2t\Delta_s}{n-1}+\frac{2t^2}{n-1}+\frac{t^2}n\right]
\le K\frac{t(\Delta_s+t)}n.
\]

保留行的误差至多 \(K(\Delta_s+t)/n\)，共 \(2m\) 个行列误差，故

\[
\|g(X)\|_2\le\sqrt{2m}\,K(\Delta_s+t)/n
\le K'(\Delta_s+t)/\sqrt n.
\]

归一化后同阶，容量和位移给出

\[
0\le\theta_{\widetilde X}\le K(\Delta_s+t)^2/n,
\quad
\|B_X-\widetilde X\|_F\le\frac K{\sqrt m}\frac{\Delta_s+t}{\sqrt n}
\le K'(\Delta_s+t)/n.
\]

因为 \(|\kappa|/m\le1/2\)，

\[
\left|m\log\frac m{m+\kappa}\right|
=m|\log(1+\kappa/m)|\le2|\kappa|,
\]

所以 \(|\theta_X|\le K[(\Delta_s+t)^2+t(\Delta_s+t)]/n\le K'(\Delta_s+t)^2/n\)。若 \(\Delta_s=t=0\)，这些量均为零。

高斯删除给 \(K(t+\tau+1)/n\)，而

\[
\tau=\frac{\sum_i s_i^2}{(n-1)^2}\le\frac{n\Delta_s^2}{(n-1)^2}\le\frac{4\Delta_s^2}n.
\]

中心核位移为 \(K(\Delta_s+t+1)/n\)，故全部高斯对数误差至多

\[
K\left[\frac{t+1}n+\frac{\Delta_s^2}{n^2}+\frac{\Delta_s+t+1}n\right]
\le K'\frac{(\Delta_s+t+1)^2}n=K'\eta.
\]

恢复上述因子后，得

\[
\frac{\operatorname{per}A[R,T]}{e^{-1}D_n(S)m!/2^m}
=e^u(1+v),\qquad |u|\le K\eta,\quad |v|\le K/n\le K\eta.
\]

选择固定 \(\delta\le1/(2K)\) 后，

\[
|e^u(1+v)-1|
\le|e^u-1|+e^{|u|}|v|
\le e^{K\delta}(|u|+|v|)
\le2Ke^{K\delta}\eta.
\]

其中 \(K,\delta,N_0\) 仅依赖 \(B_0\) 及固定的分析参数。

这里没有成对预缩放，\(C=J/(n-1)+Y\)，\(X=(n/m)C[R,T]\)，\(\widetilde X=mX/(m+\kappa)\)。所以

\[
\widetilde X=\frac n{m+\kappa}C[R,T],\qquad
\Pi_m\widetilde X\Pi_m=\frac n{m+\kappa}Z.
\]

标量偏移为

\[
\left|\frac n{m+\kappa}-1\right|
=\frac{|t-\kappa|}{|m+\kappa|}
\le\frac{4(t+|\kappa|)}n
\quad(m\ge n/2,\ |\kappa|\le m/2).
\]

于是

\[
\begin{aligned}
\|(B_X-P_m)-Z\|_F
&\le\|\Pi_m(B_X-\widetilde X)\Pi_m\|_F
+\left|\frac n{m+\kappa}-1\right|\|Z\|_F\\
&\le K\frac{\Delta_s+t}{n}+K\frac{t+|\kappa|}{n}\\
&\le K'\left[\frac{\Delta_s+t}{n}+\frac{t(\Delta_s+t)}{n^2}\right].
\end{aligned}
\]

由于 \(t/n\le1/2\)，最后一项被 \(K(\Delta_s+t)/n\) 吸收。若 \(\Delta_s=t=0\)，该位移为零。再计入 \(Y\) 与 \(S/n\) 之间的 \(O(n^{-1})\) 误差，得到统一误差界 \(O((\Delta_s+t+1)^2/n)\)。

对 \(\Delta_s,t\ge0\)，以下各项分别都不超过 \(\eta=(\Delta_s+t+1)^2/n\) 的固定倍数：

\[
\frac{t+1}n\le\eta,\quad
\frac{\Delta_s+t}n\le\eta,\quad
\frac{(\Delta_s+t)^2}n\le\eta,\quad
\frac{t(\Delta_s+t)}n\le\eta,\quad
\frac{\Delta_s^2}{n^2}\le\eta.
\]

因此，各对数误差均为 \(O(\eta)\)。由 \(|e^u-1|\le e^{|u|}|u|\)，相对误差也为 \(O(\eta)\)。 \(\square\)

\phantomsection\label{technical-spectral-cap}

\Needspace{12\baselineskip}
\section{6.1 谱半径与统一粗界}\label{ux8c31ux534aux5f84ux4e0eux7edfux4e00ux7c97ux754c}

\textbf{引理 6.1（谱上限）。} 任意 \(n\) 阶竞赛符号矩阵满足

\[
\|S\|_{\rm op}\le\cot\frac{\pi}{2n}<\frac{2n}{\pi},\qquad
\rho_n(S)\le\frac{3n-1}{n+1}<3.
\]

\textbf{证明。} \(iS\) 为 Hermitian 矩阵。对任意复向量 \(z\)，同时对 \(S,z\) 作带符号置换，使非零坐标相位落在 \([0,\pi)\) 且递增。于是 \(c_{ij}=\operatorname{Im}(\overline z_i z_j)\ge0\)（\(i<j\)）。若 \(T_n^0\) 是上三角全一的传递符号矩阵，则

\[
z^*iSz=-2\sum_{i<j}S_{ij}c_{ij}\le2\sum_{i<j}c_{ij}=z^*(-iT_n^0)z.
\]

取 Rayleigh 商得 \(\|S\|_{\rm op}\le\|T_n^0\|_{\rm op}\)。相邻行的特征方程给出 \((\lambda-1)v_i=(\lambda+1)v_{i+1}\)，几何比 \(r=(\lambda-1)/(\lambda+1)\) 满足 \(r^n=-1\)。全部特征值因而为 \(i\cot((2j-1)\pi/(2n))\)，\(1\le j\le n\)，得到所述算子界；\(n=1\) 直接验证。

令 \(x_j=\lambda_j^2/n^2\)，则 \(\sum_jx_j=(n-1)/(2n)\)。函数 \(\phi(x)=\log((1+x)/(1-x))\) 满足 \(\phi(0)=0\)，且 \(\phi'\) 在 \([0,1)\) 递增，故当 \(a,b\ge0\)、\(a+b<1\) 时，\(\phi(a)+\phi(b)\le\phi(a+b)\)。迭代即得

\[
\log\rho_n(S)=\sum_j\phi(x_j)\le\phi\left(\frac{n-1}{2n}\right),
\]

取指数得到粗界。它不使用后面的 Fang--Huang 比较。

传递矩阵的同一特征值计算给出全部偶阶主子式为一、奇阶主子式为零（偶阶正频率按互补角配对，其乘积为一）。因此

\[
\det(I+zT_n^0)=\sum_{h=0}^{\lfloor n/2\rfloor}\binom n{2h}z^{2h}
=\frac{(1+z)^n+(1-z)^n}{2}.
\]

\(\square\)

\phantomsection\label{technical-algebra}

\Needspace{12\baselineskip}
\section{A. 辅助代数与模型计算}\label{a.-ux8f85ux52a9ux4ee3ux6570ux4e0eux6a21ux578bux8ba1ux7b97}

\Needspace{12\baselineskip}
\subsection{A.1 传递符号矩阵的行列式多项式}\label{a.1-ux4f20ux9012ux7b26ux53f7ux77e9ux9635ux7684ux884cux5217ux5f0fux591aux9879ux5f0f}

引理 6.1 已由传递矩阵的特征值及偶阶主子式给出

\[
\det(I+tT_n^0)=\frac{(1+t)^n+(1-t)^n}{2}.
\]

该公式同时覆盖奇阶和偶阶，供主文的谱比较使用。

\Needspace{12\baselineskip}
\subsection{A.2 交替切换与圆环的普通同构}\label{a.2-ux4ea4ux66ffux5207ux6362ux4e0eux5706ux73afux7684ux666eux901aux540cux6784}

设\(d=2q+1\)。对传递矩阵作交替切换，\(\varepsilon_i=(-1)^{i-1}\)、\(D_{\varepsilon}=\operatorname{diag}(\varepsilon_i)\)。当\(i<j\)、\(\Delta=j-i\)时，

\[
(D_{\varepsilon}T_dD_{\varepsilon})_{ij}=(-1)^\Delta.
\]

定义圆环标签映射

\[
\phi(i)=(q+1)i\pmod d.
\]

因为\(2(q+1)=d+1\)，\(q+1\)与\(d\)互素，故\(\phi\)是置换。

若\(\Delta=2r\)，则\(1\le r\le q\)且

\[
\phi(j)-\phi(i)\equiv(q+1)2r\equiv r\pmod d.
\]

该差在\(\{1,\ldots,q\}\)，正是圆环中\(\phi(i)\to\phi(j)\)的条件；切换矩阵此时也为1。

若\(\Delta=2r+1\)，则\(0\le r\le q-1\)，且

\[
\phi(j)-\phi(i)\equiv q+1+r\pmod d.
\]

该差在\(\{q+1,\ldots,2q\}\)，方向相反；切换矩阵此时为\(-1\)。反对称性处理\(i>j\)，对角均为0，所以这是逐条边的普通同构。

主文的谱论证还用到传递矩阵的幂有相同对角元。令\(Ze_j=e_{j+1}\)（\(j<d\)）、\(Ze_d=-e_1\)。当\(i,j<d\)时，

\[
(Z^TT_dZ)_{ij}=(T_d)_{i+1,j+1}=(T_d)_{ij}.
\]

当\(i<d,j=d\)时，

\[
(Z^TT_dZ)_{i,d}=-(T_d)_{i+1,1}=1=(T_d)_{i,d}.
\]

当\(i=d,j<d\)时，

\[
(Z^TT_dZ)_{d,j}=-(T_d)_{1,j+1}=-1=(T_d)_{d,j}.
\]

最后对角元均为0。因此\(Z^TT_dZ=T_d\)，从而对每个整数\(p\ge0\)，

\[
Z^TT_d^{2p}Z=T_d^{2p}.
\]

取对角元便得到相邻对角相等，绕回时的两个负号抵消。

\Needspace{12\baselineskip}
\subsection{A.3 四阶行列式与四次矩等号}\label{a.3-ux56dbux9636ux884cux5217ux5f0fux4e0eux56dbux6b21ux77e9ux7b49ux53f7}

设 \(a,b,c,g,e,f\in\mathbb R\)，并令

\[
M=\begin{pmatrix}
0&a&b&c\\
-a&0&g&e\\
-b&-g&0&f\\
-c&-e&-f&0
\end{pmatrix}.
\]

记\(M_{1j}\)为删去第1行、第\(j\)列后的三阶行列式（不含代数余子式符号）。沿第一行展开，分别计算：

\[
\begin{aligned}
M_{12}
&=\det\begin{pmatrix}-a&g&e\\-b&0&f\\-c&-f&0\end{pmatrix}\\
&=(-a)(0-f(-f))-g((-b)0-f(-c))+e((-b)(-f)-0(-c))\\
&=-af^2-cgf+bef,\\
M_{13}
&=\det\begin{pmatrix}-a&0&e\\-b&-g&f\\-c&-e&0\end{pmatrix}\\
&=(-a)((-g)0-f(-e))-0+e((-b)(-e)-(-g)(-c))\\
&=-aef+be^2-cge,\\
M_{14}
&=\det\begin{pmatrix}-a&0&g\\-b&-g&0\\-c&-e&-f\end{pmatrix}\\
&=(-a)((-g)(-f)-0(-e))-0+g((-b)(-e)-(-g)(-c))\\
&=-agf+bge-cg^2.
\end{aligned}
\]

将这些余子式代入第一行的展开，得到

\[
\begin{aligned}
\det M
&=-aM_{12}+bM_{13}-cM_{14}\\
&=(a^2f^2+acgf-abef)
+(-abef+b^2e^2-bcge)
+(acgf-bcge+c^2g^2)\\
&=a^2f^2+b^2e^2+c^2g^2-2abef+2acgf-2bcge\\
&=(af-be+cg)^2.
\end{aligned}
\]

六个变量均为\(\pm1\)时，\(af-be+cg\)是三个\(\pm1\)之和，其值为\(\pm1\)或\(\pm3\)，故行列式为1或9。设\(D_4(S)\)为四顶点主子式等于9的个数。由行列式的列多线性展开，在每列选择单位列或\(tS\)的列，得到

\[
\det(I+tS)=\sum_{U\subseteq[d]}t^{|U|}\det S[U].
\]

奇阶反对称矩阵\(B\)满足\(\det B=\det B^T=\det(-B)=-\det B\)，所以其行列式为0。二阶主子式均为1，四阶中有\(D_4\)个9，其余为1。因此

\[
\det(I+tS)=1+e_2t^2+e_4t^4+O(t^6),
\]

\[
e_2=\binom d2,\qquad
e_4=(\binom d4-D_4)+9D_4=\binom d4+8D_4.
\]

下列展开只在\(t=0\)附近使用。取\(|t|\,\|S\|_{\rm op}<1/2\)，则每个特征值\(\mu\)都满足\(|t\mu|<1/2\)，并用在1处为0的对数级数定义\(\log(1+t\mu)\)。有限个特征值求和与幂级数交换合法，且

\[
\det(I+tS)=\prod_\mu(1+t\mu),\qquad
\operatorname{tr}S^k=\sum_\mu\mu^k.
\]

这里特征值按重数计入；实反对称标准形保证可在复数域对角化，零特征值也包括在内。

每个特征值的对数级数给

\[
\log\det(I+tS)
=\sum_{k\ge1}\frac{(-1)^{k+1}t^k}{k}\operatorname{tr}S^k
=-\frac{t^2}{2}\operatorname{tr}S^2
-\frac{t^4}{4}\operatorname{tr}S^4+O(t^6).
\]

奇次迹为0，因为\(\operatorname{tr}S^{2j+1}=\operatorname{tr}(S^{2j+1})^T=-\operatorname{tr}S^{2j+1}\)。另一方面，将\(z=e_2t^2+e_4t^4+O(t^6)\)代入\(\log(1+z)=z-z^2/2+O(z^3)\)，其中\(z^2=e_2^2t^4+O(t^6)\)、\(z^3=O(t^6)\)，得到

\[
\log\det(I+tS)=e_2t^2+(e_4-e_2^2/2)t^4+O(t^6).
\]

比较 \(t^4\) 的系数，得

\[
-\operatorname{tr}S^4/4=e_4-e_2^2/2,\qquad
\operatorname{tr}S^4=2e_2^2-4e_4
=2\left(\binom d2\right)^2-4\binom d4-32D_4(S).
\]

传递矩阵任取四个升序顶点都有\(a=b=c=g=e=f=1\)，其行列式为\((1-1+1)^2=1\)，所以\(D_4(T_d)=0\)。相减给

\[
\operatorname{tr}T_d^4-\operatorname{tr}S^4=32D_4(S).
\]

若四次矩取等，则\(D_4(S)=0\)。取\(\varepsilon_1=1,\varepsilon_j=s_{1j}\)（\(j>1\)），定义\(D_{\varepsilon}=\operatorname{diag}(\varepsilon_1,\ldots,\varepsilon_d)\)和\(S'=D_{\varepsilon} S D_{\varepsilon}\)。于是\(s'_{1j}=\varepsilon_1s_{1j}\varepsilon_j=s_{1j}^2=1\)。各主子式乘以\((\det D_{\varepsilon}[U])^2=1\)，故\(D_4\)不变。任取\(i<j<k\)且都不为1，

\[
\det S'[\{1,i,j,k\}]
=(s'_{jk}-s'_{ik}+s'_{ij})^2.
\]

若这三点形成有向三圈，则\(s'_{ij}=s'_{jk}=-s'_{ik}\)，括号为\(3s'_{ij}\)，平方为9，矛盾。因此剩余顶点无有向三圈。如果\(i\to j\)、\(j\to k\)而\(i\not\to k\)，则\(k\to i\)，会形成三圈；所以方向关系是严格全序，可按该全序排列。第一顶点指向所有剩余点，整个切换图于是为传递图。反过来带符号置换相似保持迹，故四次矩取等。\(d=3\)时切换后剩下两点，按其唯一边方向排列即可。

\phantomsection\label{technical-references}

\Needspace{12\baselineskip}
\section{参考文献}\label{ux53c2ux8003ux6587ux732e}

{[}1{]} John Irving and Mohamed Omar. Revisiting the Rédei-Berge Symmetric Functions via Matrix Algebra. The Electronic Journal of Combinatorics 32(4) (2025), Paper P4.43, 23 pp.~\href{https://doi.org/10.37236/13841}{DOI: 10.37236/13841}.

{[}2{]} Noga Alon. The Maximum Number of Hamiltonian Paths in Tournaments. Combinatorica 10(4) (1990), 319--324. \href{https://doi.org/10.1007/BF02128667}{DOI: 10.1007/BF02128667}.

{[}3{]} Yanjun Han and Jonathan Niles-Weed. Approximate independence of permutation mixtures. The Annals of Statistics, to appear（\href{https://yanjunhan2021.github.io/publication.html}{作者发表目录}）。\href{https://arxiv.org/html/2408.09341v2}{arXiv:2408.09341v2}（2024年9月9日）；正文所引引理编号对应这一版本。

{[}4{]} Bo Deng, Xueliang Li, Bryan Shader and Wasin So. On the Maximum Skew Spectral Radius and Minimum Skew Energy of Tournaments. Linear and Multilinear Algebra 66(7) (2018), 1434--1441. \href{https://doi.org/10.1080/03081087.2017.1357676}{DOI: 10.1080/03081087.2017.1357676}.

{[}5{]} Peter McCullagh. An asymptotic approximation for the permanent of a doubly stochastic matrix. Journal of Statistical Computation and Simulation 84(2) (2014), 404--414. \href{https://doi.org/10.1080/00949655.2012.712122}{DOI: 10.1080/00949655.2012.712122}. \href{https://arxiv.org/abs/1205.5723}{arXiv:1205.5723}.

{[}6{]} N. C. Wormald. Tournaments with many Hamilton cycles. arXiv:2610.16245v1 (2026). \href{https://arxiv.org/html/2610.16245v1}{Version}.

{[}7{]} Ehud Friedgut and Jeff Kahn. On the Number of Hamiltonian Cycles in a Tournament. Combinatorics, Probability and Computing 14(5--6) (2005), 769--781. \href{https://doi.org/10.1017/S0963548305006863}{DOI: 10.1017/S0963548305006863}.

{[}8{]} Eric Li. The Godsil--McKay Asymptotic for Latin Rectangles in the Sublinear Range of Erdős Problem 725. \href{https://arxiv.org/html/2608.01671v1}{arXiv:2608.01671v1} (2026).

{[}9{]} Ilan Adler, Noga Alon and Sheldon M. Ross. On the Maximum Number of Hamiltonian Paths in Tournaments. Random Structures \& Algorithms 18(3) (2001), 291--296. \href{https://doi.org/10.1002/rsa.1010}{DOI: 10.1002/rsa.1010}.

{[}10{]} Tibor Szele. Kombinatorikai vizsgálatok az irányított teljes gráffal kapcsolatban. Matematikai és Fizikai Lapok 50 (1943), 223--256. \href{https://real-j.mtak.hu/7300/}{Original volume archive}.

{[}11{]} Yijia Fang and Hao Huang. Maximizing directed cycles in tournaments. arXiv:2608.06238v1 (2026). \href{https://arxiv.org/html/2608.06238v1\#S3}{arXiv version}.

{[}12{]} Alexander Divoux and Shouda Wang. The maximum number of Hamiltonian cycles in tournaments. arXiv:2610.23755v1 (2026). \href{https://arxiv.org/html/2610.23755v1}{arXiv version}.

{[}13{]} Deep Bhattacharjee. \emph{Two conjectures of Wormald on Hamilton cycles in tournaments}. Hexagon preprint 2610.00460v2 (2026). \href{https://hexagonmath.org/2610.00460v2}{Version}.

{[}14{]} Xingchen Liu. \emph{Constant-factor bounds for Hamiltonian paths in tournaments}, v1.0 (8 October 2026). \href{https://doi.org/10.5281/zenodo.23233166}{Version DOI: 10.5281/zenodo.23233166}. Later archive: v1.2, \href{https://doi.org/10.5281/zenodo.23288896}{10.5281/zenodo.23288896}.

{[}15{]} Xingchen Liu and Xiangyu Ye. Wormald tournament project, manuscript revision \nolinkurl{2fd6a2c30260b826feef71049d7d593527fd7a42} (19 October 2026). \href{https://github.com/LStar404/wormald-tournament-conjectures/blob/2fd6a2c30260b826feef71049d7d593527fd7a42/docs/paper.zh.md}{Fixed manuscript}.

{[}16{]} Deep Bhattacharjee. \emph{Two conjectures of Wormald on Hamilton cycles in tournaments}, release v1.0.0 (14 October 2026). \href{https://github.com/creelie/Breakthrough-problems-1/releases/tag/v1.0.0}{Release and source archive}.

{[}17{]} Liu--Ye. \emph{Tournament Hamiltonian counts}: source and formal-coverage snapshot \nolinkurl{c678e9ee24fb304cad755a27c99f511c8640f75e}. \href{https://github.com/LStar404/tournament-hamiltonian-counts-20261028/tree/c678e9ee24fb304cad755a27c99f511c8640f75e}{Fixed repository revision}.

\end{document}
